\documentclass[12pt,a4paper]{article}
\usepackage[UTF8]{ctex}
\usepackage{amsmath,amssymb,amsthm}
\usepackage{geometry}
\usepackage{graphicx}
\usepackage{hyperref}

\geometry{margin=2.5cm}

\title{第8章：开链动力学}
\author{Modern Robotics}
\date{}

\begin{document}

\maketitle

Chapter 8

Dynamics of Open Chains

In this chapter we study once again the motions of open-chain robots, but this
time taking into account the forces and torques that cause them; this is the
subject of robot dynamics. The associated dynamic equations – also referred
to as the equations of motion – are a set of second-order differential equations
of the form
\begin{equation}
τ = M (θ)θ̈ + h(θ, θ̇),
\label{eq:8.1}
\end{equation}
where θ ∈ Rn is the vector of joint variables, τ ∈ Rn is the vector of joint forces
and torques, M (θ) ∈ Rn×n is a symmetric positive-definite mass matrix, and
h(θ, θ̇) ∈ Rn are forces that lump together centripetal, Coriolis, gravity, and
friction terms that depend on θ and θ̇. One should not be deceived by the
apparent simplicity of these equations; even for “simple” open chains, e.g., those
with joint axes that are either orthogonal or parallel to each other, M (θ) and
h(θ, θ̇) can be extraordinarily complex.
    Just as a distinction was made between a robot’s forward and inverse kine-
matics, it is also customary to distinguish between a robot’s forward and in-
verse dynamics. The forward problem is the problem of determining the
robot’s acceleration θ̈ given the state (θ, θ̇) and the joint forces and torques,
                                                     
\begin{equation}
θ̈ = M −1 (θ) τ − h(θ, θ̇) ,
\label{eq:8.2}
\end{equation}
and the inverse problem is finding the joint forces and torques τ corresponding
to the robot’s state and a desired acceleration, i.e., Equation (8.1).
    A robot’s dynamic equations are typically derived in one of two ways: by a
direct application of Newton’s and Euler’s dynamic equations for a rigid body
(often called the Newton–Euler formulation) or by the Lagrangian dy-
namics formulation derived from the kinetic and potential energy of the robot.

                                        271
272                                                                 8.1. Lagrangian Formulation

The Lagrangian formalism is conceptually elegant and quite effective for robots
with simple structures, e.g., with three or fewer degrees of freedom. The calcula-
tions can quickly become cumbersome for robots with more degrees of freedom,
however. For general open chains, the Newton–Euler formulation leads to effi-
cient recursive algorithms for both the inverse and forward dynamics that can
also be assembled into closed-form analytic expressions for, e.g., the mass matrix
M (θ) and the other terms in the dynamics equation (8.1). The Newton–Euler
formulation also takes advantage of tools we have already developed in this
book.
    In this chapter we study both the Lagrangian and Newton–Euler dynamics
formulations for an open-chain robot. While we usually express the dynamics
in terms of the joint space variables θ, it is sometimes convenient to express it
in terms of the configuration, twist, and rate of change of the twist of the end-
effector. This is the task-space dynamics, studied in Section 8.6. Sometimes
robots are subject to a set of constraints on their motion, such as when the
robot makes contact with a rigid environment. This leads to a formulation of
the constrained dynamics (Section 8.7), whereby the space of joint torques and
forces is divided into a subspace that causes motion of the robot and a subspace
that causes forces against the constraints. The URDF file format for specifying
robot inertial properties is described in Section 8.8. Finally, some practical
issues that arise in the derivation of robot dynamics, such as the effect of motor
gearing and friction, are described in Section 8.9.

\section{Lagrangian Formulation}
\subsection{Basic Concepts and Motivating Examples}
The first step in the Lagrangian formulation of dynamics is to choose a set of
independent coordinates q ∈ Rn that describes the system’s configuration. The
coordinates q are called generalized coordinates. Once generalized coordi-
nates have been chosen, these then define the generalized forces f ∈ Rn . The
forces f and the coordinate rates q̇ are dual to each other in the sense that the
inner product f T q̇ corresponds to power. A Lagrangian function L(q, q̇) is then
defined as the overall system’s kinetic energy K(q, q̇) minus the potential energy
P(q),
                             L(q, q̇) = K(q, q̇) − P(q).
The equations of motion can now be expressed in terms of the Lagrangian as
follows:
                                 d ∂L ∂L
\begin{equation}
f=         −    ,
\label{eq:8.3}
\end{equation}
                                dt ∂ q̇   ∂q


Chapter 8. Dynamics of Open Chains                                                                  273

                                            m2
                               L2
                                            θ2
                ŷ
                                       m1                                    θ2 = π/2
                        L1
                             θ1                      g

                                  x̂                              θ1 = 0

     \begin{figure}[h]
\centering
\includegraphics{figure_8_1.pdf}
\caption{(Left) A 2R open chain under gravity. (Right) At θ = (0, π/2).}
\label{fig:8_1}
\end{figure}

These equations are also referred to as the Euler–Lagrange equations with
external forces.1 The derivation can be found in dynamics texts.
   We illustrate the Lagrangian dynamics formulation through two examples.
In the first example, consider a particle of mass m constrained to move on
a vertical line. The particle’s configuration space is this vertical line, and a
natural choice for a generalized coordinate is the height of the particle, which
we denote by the scalar variable x ∈ R. Suppose that the gravitational force mg
acts downward, and an external force f is applied upward. By Newton’s second
law, the equation of motion for the particle is
\begin{equation}
f − mg = mẍ.
\label{eq:8.4}
\end{equation}
We now apply the Lagrangian formalism to derive the same result. The kinetic
energy is mẋ2 /2, the potential energy is mgx, and the Lagrangian is
                                                                 1 2
\begin{equation}
L(x, ẋ) = K(x, ẋ) − P(x) =               mẋ − mgx.
\label{eq:8.5}
\end{equation}
                                                                 2
The equation of motion is then given by
                                            d ∂L ∂L
\begin{equation}
f=                −    = mẍ + mg,
\label{eq:8.6}
\end{equation}
                                            dt ∂ ẋ   ∂x
which matches Equation (8.4).
    We now derive the dynamic equations for a planar 2R open chain moving
in the presence of gravity (Figure 8.1). The chain moves in the x̂–ŷ-plane, with
gravity g acting in the −ŷ-direction. Before the dynamics can be derived, the
mass and inertial properties of all the links must be specified. To keep things
simple the two links are modeled as point masses m1 and m2 concentrated at
   1 The external force f is zero in the standard form of the Euler–Lagrange equations.


274                                                                 8.1. Lagrangian Formulation

the ends of each link. The position and velocity of the link-1 mass are then
given by
                                             
                         x1           L1 cos θ1
                               =                  ,
                         y1           L1 sin θ1
                                               
                         ẋ1          −L1 sin θ1
                               =                    θ̇1 ,
                         ẏ1           L1 cos θ1
while those of the link-2 mass are given by
                                              
      x2            L1 cos θ1 + L2 cos(θ1 + θ2 )
             =                                     ,
      y2            L1 sin θ1 + L2 sin(θ1 + θ2 )
                                                                        
      ẋ2           −L1 sin θ1 − L2 sin(θ1 + θ2 ) −L2 sin(θ1 + θ2 )      θ̇1
             =                                                                 .
      ẏ2            L1 cos θ1 + L2 cos(θ1 + θ2 )    L2 cos(θ1 + θ2 )    θ̇2
We choose the joint coordinates θ = (θ1 , θ2 ) as the generalized coordinates.
The generalized forces τ = (τ1 , τ2 ) then correspond to joint torques (since τ T θ̇
corresponds to power). The Lagrangian L(θ, θ̇) is of the form
                                                  2
                                                  X
\begin{equation}
L(θ, θ̇) =     (Ki − Pi ),
\label{eq:8.7}
\end{equation}
                                                  i=1

where the link kinetic energy terms K1 and K2 are
            1                    1
K1    =       m1 (ẋ21 + ẏ12 ) = m1 L21 θ̇12
            2                    2
            1        2      2
K2    =       m2 (ẋ2 + ẏ2 )
            2
            1                                                                              
      =       m2 (L21 + 2L1 L2 cos θ2 + L22 )θ̇12 + 2(L22 + L1 L2 cos θ2 )θ̇1 θ̇2 + L22 θ̇22 ,
            2
and the link potential energy terms P1 and P2 are
                   P1    = m1 gy1 = m1 gL1 sin θ1 ,
                   P2    = m2 gy2 = m2 g(L1 sin θ1 + L2 sin(θ1 + θ2 )).
The Euler–Lagrange equations (8.3) for this example are of the form
                                      d ∂L       ∂L
\begin{equation}
τi =            −     ,        i = 1, 2.
\label{eq:8.8}
\end{equation}
                                      dt ∂ θ˙i   ∂θi
The dynamic equations for the 2R planar chain follow from explicit evaluation
of the right-hand side of (8.8) (we omit the detailed calculations, which are
straightforward but tedious):


Chapter 8. Dynamics of Open Chains                                                                    275

                                                                              
       τ1   =       m1 L21 + m2 (L21 + 2L1 L2 cos θ2 + L22 ) θ̈1               
                                                                               
                                                                               
                                                                               
                                                                             2 
              + m2 (L1 L2 cos θ2 + L2 )θ̈2 − m2 L1 L2 sin θ2 (2θ̇1 θ̇2 + θ̇2 ) 
                                       2
                                                                               
                                                                               
\begin{equation}
+ (m1 + m2 )L1 g cos θ1 + m2 gL2 cos(θ1 + θ2 ),
\label{eq:8.9}
\end{equation}
                                                                               
                                                                               
       τ2   = m2 (L1 L2 cos θ2 + L22 )θ̈1 + m2 L22 θ̈2 + m2 L1 L2 θ̇12 sin θ2 
                                                                               
                                                                               
                                                                               
                                                                               
              + m2 gL2 cos(θ1 + θ2 ).

    We can gather terms together into an equation of the form
\begin{equation}
τ = M (θ)θ̈ + c(θ, θ̇) + g(θ),
\label{eq:8.10}
\end{equation}
                                                |      {z     }
                                                         h(θ,θ̇)

with
                                                                                                 
                m1 L21 + m2 (L21 + 2L1 L2 cos θ2 + L22 ) m2 (L1 L2 cos θ2 + L22 )
    M (θ) =                                                                                           ,
                        m2 (L1 L2 cos θ2 + L22 )                  m2 L22
                                                   
                −m2 L1 L2 sin θ2 (2θ̇1 θ̇2 + θ̇22 )
   c(θ, θ̇) =                                         ,
                       m2 L1 L2 θ̇12 sin θ2
                                                             
                (m1 + m2 )L1 g cos θ1 + m2 gL2 cos(θ1 + θ2 )
     g(θ) =                                                     ,
                             m2 gL2 cos(θ1 + θ2 )

where M (θ) is the symmetric positive-definite mass matrix, c(θ, θ̇) is the vector
containing the Coriolis and centripetal torques, and g(θ) is the vector containing
the gravitational torques. These reveal that the equations of motion are linear
in θ̈, quadratic in θ̇, and trigonometric in θ. This is true in general for serial
chains containing revolute joints, not just for the 2R robot.
    The M (θ)θ̈ + c(θ, θ̇) terms in Equation (8.10) could have been derived by
writing fi = mi ai for each point mass, where the accelerations ai are written
in terms of θ, by differentiating the expressions for (ẋ1 , ẏ1 ) and (ẋ2 , ẏ2 ) given
above:
                                                                   
           fx1               ẍ1                −L1 θ̇12 c1 − L1 θ̈1 s1
\begin{equation}
f1 =  fy1  = m1  ÿ1  = m1  −L1 θ̇12 s1 + L1 θ̈1 c1  ,
\label{eq:8.11}
\end{equation}
           fz1               z̈1                           0
                                                                                  
               −L1 θ̇12 c1 − L2 (θ̇1 + θ̇2 )2 c12 − L1 θ̈1 s1 − L2 (θ̈1 + θ̈2 )s12
\begin{equation}
f2 = m2  −L1 θ̇12 s1 − L2 (θ̇1 + θ̇2 )2 s12 + L1 θ̈1 c1 + L2 (θ̈1 + θ̈2 )c12  ,
\label{eq:8.12}
\end{equation}
                                                0


276                                                                 8.1. Lagrangian Formulation

where s12 indicates sin(θ1 + θ2 ), etc. Defining r11 as the vector from joint 1 to
m1 , r12 as the vector from joint 1 to m2 , and r22 as the vector from joint 2 to
m2 , the moments in world-aligned frames {i} attached to joints 1 and 2 can be
expressed as m1 = r11 × f1 + r12 × f2 and m2 = r22 × f2 . (Note that joint 1
must provide torques to move both m1 and m2 , but joint 2 only needs to provide
torque to move m2 .) The joint torques τ1 and τ2 are just the third elements of
m1 and m2 , i.e., the moments about the ẑi axes out of the page, respectively.
     In (x, y) coordinates, the accelerations of the masses are written simply as
second time-derivatives of the coordinates, e.g., (ẍ2 , ÿ2 ). This is because the
x̂-ŷ frame is an inertial frame. The joint coordinates (θ1 , θ2 ) are not in an
inertial frame, however, so accelerations are expressed as a sum of terms that
are linear in the second derivatives of joint variables, θ̈, and quadratic of the
first derivatives of joint variables, θ̇T θ̇, as seen in Equations (8.11) and (8.12).
Quadratic terms containing θ̇i2 are called centripetal terms, and quadratic
terms containing θ̇i θ̇j , i 6= j, are called Coriolis terms. In other words, θ̈ = 0
does not mean zero acceleration of the masses, due to the centripetal and Coriolis
terms.
     To better understand the centripetal and Coriolis terms, consider the arm
at the configuration (θ1 , θ2 ) = (0, π/2), i.e., cos θ1 = sin(θ1 + θ2 ) = 1, sin θ1 =
cos(θ1 + θ2 ) = 0. Assuming θ̈ = 0, the acceleration (ẍ2 , ÿ2 ) of m2 from Equa-
tion (8.12) can be written
                                                                 
                    ẍ2               −L1 θ̇12              0
                            =                         +                 .
                     ÿ2          −L2 θ̇12 − L2 θ̇22     −2L2 θ̇1 θ̇2
                               |          {z         } |    {z        }
                                      centripetal terms          Coriolis terms

Figure 8.2 shows the centripetal acceleration acent1 = (−L1 θ̇12 , −L2 θ̇12 ) when
θ̇2 = 0, the centripetal acceleration acent2 = (0, −L2 θ̇22 ) when θ̇1 = 0, and the
Coriolis acceleration acor = (0, −2L2 θ̇1 θ̇2 ) when both θ̇1 and θ̇2 are positive.
As illustrated in Figure 8.2, each centripetal acceleration acenti pulls m2 toward
joint i to keep m2 rotating about the center of the circle defined by joint i.2
Therefore acenti creates zero torque about joint i. The Coriolis acceleration acor
in this example passes through joint 2, so it creates zero torque about joint 2
but it creates negative torque about joint 1; the torque about joint 1 is negative
because m2 gets closer to joint 1 (due to joint 2’s motion). Therefore the inertia
due to m2 about the ẑ1 -axis is dropping, meaning that the positive momentum
about joint 1 drops while joint 1’s speed theta ˙ 1 is constant. Therefore joint 1
must apply a negative torque, since torque is defined as the rate of change of
    2 Without this centripetal acceleration, and therefore centripetal force, the mass m would
                                                                                        2
fly off along a tangent to the circle.


Chapter 8. Dynamics of Open Chains                                                                  277

           acent1                                                             acent1
                                              acent2
                                                                          acent2

                                                                             acor

\begin{figure}[h]
\centering
\includegraphics{figure_8_2.pdf}
\caption{Accelerations of m2 when θ = (0, π/2) and θ̈ = 0. (Left) The centripetal}
\label{fig:8_2}
\end{figure}
acceleration acent1 = (−L1 θ̇12 , −L2 θ̇12 ) of m2 when θ̇2 = 0. (Middle) The centripetal
acceleration acent2 = (0, −L2 θ̇22 ) of m2 when θ̇1 = 0. (Right) When both joints are
rotating with θ̇i > 0, the acceleration is the vector sum of acent1 , acent2 , and the
Coriolis acceleration acor = (0, −2L2 θ̇1 θ̇2 ).

angular momentum. Otherwise θ̇1 would increase as m2 gets closer to joint 1,
just as a skater’s rotation speed increases as she pulls in her outstretched arms
while doing a spin.

\subsection{General Formulation}
We now describe the Lagrangian dynamics formulation for general n-link open
chains. The first step is to select a set of generalized coordinates θ ∈ Rn for
the configuration space of the system. For open chains all of whose joints are
actuated, it is convenient and always possible to choose θ to be the vector of the
joint values. The generalized forces will be denoted τ ∈ Rn . If θi is a revolute
joint then τi will correspond to a torque, while if θi is a prismatic joint then τi
will correspond to a force.
    Once θ has been chosen and the generalized forces τ identified, the next step
is to formulate the Lagrangian L(θ, θ̇) as follows:
\begin{equation}
L(θ, θ̇) = K(θ, θ̇) − P(θ),
\label{eq:8.13}
\end{equation}
where K(θ, θ̇) is the kinetic energy and P(θ) is the potential energy of the overall
system. For rigid-link robots the kinetic energy can always be written in the
form
                                n   n
                             1 XX                      1
\begin{equation}
K(θ) =             mij (θ)θ̇i θ̇j = θ̇T M (θ)θ̇,
\label{eq:8.14}
\end{equation}
                             2 i=1 j=1                 2
where mij (θ) is the (i, j)th element of the n × n mass matrix M (θ); a con-
structive proof of this assertion is provided when we examine the Newton–Euler
formulation.


278                                                                     8.1. Lagrangian Formulation

   The dynamic equations are analytically obtained by evaluating the right-
hand side of
                        d ∂L       ∂L
\begin{equation}
τi =           −     ,  i = 1, . . . , n.
\label{eq:8.15}
\end{equation}
                        dt ∂ θ̇i   ∂θi
With the kinetic energy expressed as in Equation (8.14), the dynamics can be
written explicitly as
             n
             X                    n X
                                  X n
                                                                ∂P
\begin{equation}
τi =         mij (θ)θ̈j +             Γijk (θ)θ̇j θ̇k +       ,       i = 1, . . . , n,
\label{eq:8.16}
\end{equation}
             j=1                  j=1 k=1
                                                                ∂θi

where the Γijk (θ), known as the Christoffel symbols of the first kind, are
defined as follows:
                                                     
                                1 ∂mij   ∂mik    ∂mjk
\begin{equation}
Γijk (θ) =        +       −        .
\label{eq:8.17}
\end{equation}
                                2 ∂θk     ∂θj     ∂θi

This shows that the Christoffel symbols, which generate the Coriolis and cen-
tripetal terms c(θ, θ̇), are derived from the mass matrix M (θ).
    As we have already seen, the equations (8.16) are often gathered together in
the form

                τ = M (θ)θ̈ + c(θ, θ̇) + g(θ)            or         M (θ)θ̈ + h(θ, θ̇),

where g(θ) is simply ∂P/∂θ.
    We can see explicitly that the Coriolis and centripetal terms are quadratic
in the velocity by using the form
\begin{equation}
τ = M (θ)θ̈ + θ̇T Γ(θ)θ̇ + g(θ),
\label{eq:8.18}
\end{equation}
where Γ(θ) is an n×n×n matrix and the product θ̇T Γ(θ)θ̇ should be interpreted
as follows:                          T            
                                      θ̇ Γ1 (θ)θ̇
                                     θ̇T Γ2 (θ)θ̇ 
                                                  
                       θ̇T Γ(θ)θ̇ =       ..      ,
                                           .      
                                                     θ̇T Γn (θ)θ̇
where Γi (θ) is an n × n matrix with (j, k)th entry Γijk .
  It is also common to see the dynamics written as

                                  τ = M (θ)θ̈ + C(θ, θ̇)θ̇ + g(θ),


Chapter 8. Dynamics of Open Chains                                                                  279

where C(θ, θ̇) ∈ Rn×n is called the Coriolis matrix, with (i, j)th entry
                                                    n
                                                    X
\begin{equation}
cij (θ, θ̇) =         Γijk (θ)θ̇k .
\label{eq:8.19}
\end{equation}
                                                    k=1

The Coriolis matrix is used to prove the following passivity property (Propo-
sition 8.1), which can be used to prove the stability of certain robot control
laws, as we will see in Section 11.4.2.2.
Proposition 8.1. The matrix Ṁ (θ) − 2C(θ, θ̇) ∈ Rn×n is skew symmetric,
where M (θ) ∈ Rn×n is the mass matrix, Ṁ (θ) its time derivative, and C(θ, θ̇) ∈
Rn×n is the Coriolis matrix as defined in Equation (8.19).
Proof. The (i, j)th component of Ṁ − 2C is
                                       n
                                       X ∂mij                ∂mij       ∂mik       ∂mkj
       ṁij (θ) − 2cij (θ, θ̇)    =                  θ̇k −        θ̇k −      θ̇k +      θ̇k
                                              ∂θk             ∂θk        ∂θj        ∂θi
                                       k=1
                                       Xn
                                             ∂mkj       ∂mik
                                  =               θ̇k −      θ̇k .
                                              ∂θi        ∂θj
                                       k=1

By switching the indices i and j, it can be seen that

                       ṁji (θ) − 2cji (θ, θ̇) = −(ṁij (θ) − 2cij (θ, θ̇)),

thus proving that (Ṁ − 2C)T = −(Ṁ − 2C) as claimed.

\subsection{Understanding the Mass Matrix}
The kinetic energy 12 θ̇T M (θ)θ̇ is a generalization of the familiar expression
1    T
2 mv v for a point mass. The fact that the mass matrix M (θ) is positive
definite, meaning that θ̇T M (θ)θ̇ > 0 for all θ̇ 6= 0, is a generalization of the fact
that the mass of a point mass is always positive, m > 0. In both cases, if the
velocity is nonzero, the kinetic energy must be positive.
    On the one hand, for a point mass with dynamics expressed in Cartesian
coordinates as f = mẍ, the mass is independent of the direction of acceleration,
and the acceleration ẍ is always “parallel” to the force, in the sense that ẍ
is a scalar multiple of f . A mass matrix M (θ), on the other hand, presents a
different effective mass in different acceleration directions, and θ̈ is not generally
a scalar multiple of τ even when θ̇ = 0. To visualize the direction dependence
of the effective mass, we can map a unit ball of joint accelerations {θ̈ | θ̈T θ̈ = 1}
through the mass matrix M (θ) to generate a joint force–torque ellipsoid when


280                                                                       8.1. Lagrangian Formulation

           θ̈2                                                              τ2
                                                           0.5 −0.5
                                        M −1 (θ) =
                                                           −0.5 1.5

                         θ1 = 0◦
                         θ2 = 90◦                           
                                                       3 1
                                         M (θ) =
                                                       1 1                                      τ1
                                  θ̈1

           θ̈2
                                                                              τ2

                                                                         
                                                            0.8  −0.11
                                        M −1 (θ) =
                                                           −0.11 1.01

                                  θ̈1
                                                                                                τ1
                                                                 
                                                      1.27    0.13
                                        M (θ) =
                         θ1 = 0◦                      0.13      1
                         θ2 = 150◦

\begin{figure}[h]
\centering
\includegraphics{figure_8_3.pdf}
\caption{(Bold lines) A unit ball of accelerations in θ̈ maps through the mass}
\label{fig:8_3}
\end{figure}
matrix M (θ) to a torque ellipsoid that depends on the configuration of the 2R arm.
These torque ellipsoids may be interpreted as mass ellipsoids. The mapping is shown
for two arm configurations: (0◦ , 90◦ ) and (0◦ , 150◦ ). (Dotted lines) A unit ball in τ
maps through M −1 (θ) to an acceleration ellipsoid.

the mechanism is at rest (θ̇ = 0). An example is shown in Figure 8.3 for
the 2R arm of Figure 8.1, with L1 = L2 = m1 = m2 = 1, at two different
joint configurations: (θ1 , θ2 ) = (0◦ , 90◦ ) and (θ1 , θ2 ) = (0◦ , 150◦ ). The torque
ellipsoid can be interpreted as a direction-dependent mass ellipsoid: the same
joint acceleration magnitude kθ̈k requires different joint torque magnitudes kτ k
depending on the acceleration direction. The directions of the principal axes
of the mass ellipsoid are given by the eigenvectors vi of M (θ) and the lengths
of the principal semi-axes are given by the corresponding eigenvalues λi . The
acceleration θ̈ is only a scalar multiple of τ when τ is along a principal axis of
the ellipsoid.
    It is easier to visualize the mass matrix if it is represented as an effective mass


Chapter 8. Dynamics of Open Chains                                                                  281

of the end-effector, since it is possible to feel this mass directly by grabbing and
moving the end-effector. If you grabbed the endpoint of the 2R robot, depending
on the direction you applied force to it, how massy would it feel? Let us denote
the effective mass matrix at the end-effector as Λ(θ), and the velocity of the
end-effector as V = (ẋ, ẏ). We know that the kinetic energy of the robot must
be the same regardless of the coordinates we use, so
                                   1 T           1
\begin{equation}
θ̇ M (θ)θ̇ = V T Λ(θ)V.
\label{eq:8.20}
\end{equation}
                                   2             2
\begin{equation}
Assuming the Jacobian J(θ) satisfying V = J(θ)θ̇ is invertible, Equation
\label{eq:8.20}
\end{equation}
can be rewritten as follows:

                                 V T ΛV = (J −1 V )T M (J −1 V )
                                           = V T (J −T M J −1 )V.

In other words, the end-effector mass matrix is
\begin{equation}
Λ(θ) = J −T (θ)M (θ)J −1 (θ).
\label{eq:8.21}
\end{equation}
Figure 8.4 shows the end-effector mass ellipsoids, with principal-axis directions
given by the eigenvectors of Λ(θ) and principal semi-axis lengths given by its
eigenvalues, for the same two 2R robot configurations as in Figure 8.3. The
endpoint acceleration (ẍ, ÿ) is a scalar multiple of the force (fx , fy ) applied at
the endpoint only if the force is along a principal axis of the ellipsoid. Unless
Λ(θ) is of the form cI, where c > 0 is a scalar and I is the identity matrix, the
mass at the endpoint feels different from a point mass.
    The change in apparent endpoint mass as a function of the configuration of
the robot is an issue for robots used as haptic displays. One way to reduce the
sensation of a changing mass to the user is to make the mass of the links as
small as possible.
    Note that the ellipsoidal interpretations of the relationship between forces
and accelerations defined here are only relevant at zero velocity, where there are
no Coriolis or centripetal terms.

\subsection{Lagrangian Dynamics vs. Newton–Euler Dynamics}
In the rest of this chapter, we focus on the Newton–Euler recursive method
for calculating robot dynamics. Using the tools we have developed so far, the
Newton–Euler formulation allows computationally efficient computer implemen-
tation, particularly for robots with many degrees of freedom, without the need


282                                                                            8.1. Lagrangian Formulation

                                                                                   fy
              ÿ
                                                                          
                                                                  1 0
                                       Λ−1 (θ) =
                                                                  0 0.5

                                                                                               fx
                                                                  
                                                          1       0
                                ẍ     Λ(θ) =
                                                          0       2                        θ1 = 0◦
                                                                                           θ2 = 90◦

              ÿ
                                                                                   fy      θ1 = 0◦
                                                                      
                               Λ(θ) =
                                       4    −1.73                                          θ2 = 150◦
                                      −1.73   2
                                                                                               fx
                                ẍ
                                                              
                                                0.4       0.35
                               Λ−1 (θ) =
                                               0.35        0.8

\begin{figure}[h]
\centering
\includegraphics{figure_8_4.pdf}
\caption{(Bold lines) A unit ball of accelerations in (ẍ, ÿ) maps through the}
\label{fig:8_4}
\end{figure}
end-effector mass matrix Λ(θ) to an end-effector force ellipsoid that depends on the
configuration of the 2R arm. For the configuration (θ1 , θ2 ) = (0◦ , 90◦ ), a force in the
fy -direction exactly feels both masses m1 and m2 , while a force in the fx -direction
feels only m2 . (Dotted lines) A unit ball in f maps through Λ−1 (θ) to an acceleration
ellipsoid. The × symbols for (θ1 , θ2 ) = (0◦ , 150◦ ) indicate an example endpoint force
(fx , fy ) = (1, 0) and its corresponding acceleration (ẍ, ÿ) = (0.4, 0.35), showing that
the force and acceleration at the endpoint are not aligned.

for differentiation. The resulting equations of motion are, and must be, identical
with those derived using the energy-based Lagrangian method.
    The Newton–Euler method builds on the dynamics of a single rigid body, so
we begin there.


Chapter 8. Dynamics of Open Chains                                                                  283

\section{Dynamics of a Single Rigid Body}
\subsection{Classical Formulation}
Consider a rigid body consisting of a number of rigidly connected   P point masses,
where point mass i has mass mi and the total mass is m = i mi . Let ri =
(xi , yi , zi ) be the fixed location of mass i in a body frame {b}, where the origin
of this frame is the unique point such that
                                        X
                                           mi ri = 0.
                                             i

This point is known as the center of mass. If some other point happens to
be inconveniently chosen as thePorigin, then the frame {b} should be moved
to the center of mass at (1/m) i mi ri (in the inconvenient frame) and the ri
recalculated in the center-of-mass frame.
    Now assume that the body is moving with a body twist Vb = (ωb , vb ), and
let pi (t) be the time-varying position of mi , initially located at ri , in the inertial
frame {b}. Then

                          ṗi = vb + ωb × pi ,
                                       d                  d
                          p̈i = v̇b + ωb × pi + ωb × pi
                                      dt                 dt
                              = v̇b + ω̇b × pi + ωb × (vb + ωb × pi ).

Substituting ri for pi on the right-hand side and using our skew-symmetric
notation (see Equation (3.30)), we get

                              p̈i = v̇b + [ω̇b ]ri + [ωb ]vb + [ωb ]2 ri .

Taking as a given that fi = mi p̈i for a point mass, the force acting on mi is

                           fi = mi (v̇b + [ω̇b ]ri + [ωb ]vb + [ωb ]2 ri ),

which implies a moment
                                            mi = [ri ]fi .
The total force and moment acting on the body is expressed as the wrench Fb :
                                       P            
                                  mb             m
                          Fb =          = Pi i .
                                  fb            i fi
                                                                       P
   To simplify Pthe expressions for fb and mb , keep in mind that i mi ri = 0
(and therefore i mi [ri ] = 0) and, for a, b ∈ R3 , [a] = −[a]T , [a]b = −[b]a, and


284                                                           8.2. Dynamics of a Single Rigid Body

[a][b] = ([b][a])T . Focusing on the linear dynamics,
                    X
              fb =     mi (v̇b + [ω̇b ]ri + [ωb ]vb + [ωb ]2 ri )
                        i

                      X                               X        *0 X               :0
                  =         mi (v̇b + [ωb ]vb ) −          i i b+
                                                            [r ]ω̇     i [r
                                                                          i ][ωb ]ωb
                                                                                    
                                                          m          m
                                                                                
                                                                     
                        i                             i          i
                      X
                  =         mi (v̇b + [ωb ]vb )
                        i
\begin{equation}
= m(v̇b + [ωb ]vb ).
\label{eq:8.22}
\end{equation}
The velocity product term m[ωb ]vb arises from the fact that, for ωb 6= 0, a
constant vb 6= 0 corresponds to a changing linear velocity in an inertial frame.
   Now focusing on the rotational dynamics,
                       X
             mb =         mi [ri ](v̇b + [ω̇b ]ri + [ωb ]vb + [ωb ]2 ri )
                               i

                              X       *0 X
                                                                 :0
                        =        i i b+
                                m[r ]v̇           m i [ri ][ωb ]vb
                                                                 
                                   
                                                         
                             i X
                                            i
                             +    mi [ri ]([ω̇b ]ri + [ωb ]2 ri )
                                   i
                              X                                               
                        =           mi −[ri ]2 ω̇b − [ri ]T [ωb ]T [ri ]ωb
                               i
                              X                                       
                        =           mi −[ri ]2 ω̇b − [ωb ][ri ]2 ωb
                               i
                                                      !                                  !
                                    X                                     X
                                                  2                                  2
                        =       −          mi [ri ]       ω̇b + [ωb ] −       mi [ri ]       ωb
                                       i                                  i
\begin{equation}
= Ib ω̇b + [ωb ]Ib ωb ,
\label{eq:8.23}
\end{equation}
                    P
where Ib = − i mi [ri ]2 ∈ R3×3 is the body’s rotational inertia matrix.
Equation (8.23) is known as Euler’s equation for a rotating rigid body.
    In Equation (8.23), note the presence of a term linear in the angular accel-
eration, Ib ω̇b , and a term quadratic in the angular velocities, [ωb ]Ib ωb , just as
we saw for the mechanisms in Section 8.1. Also, Ib is symmetric and positive
definite, just like the mass matrix for a mechanism, and the rotational kinetic
energy is given by the quadratic
                                                      1 T
                                              K=       ω I b ωb .
                                                      2 b


Chapter 8. Dynamics of Open Chains                                                                  285

One difference is that Ib is constant whereas the mass matrix M (θ) changes
with the configuration of the mechanism.
   Writing out the individual entries of Ib , we get
                  P        2    2
                                         P               P              
                        i (yi + zi ) P
                       mP              − mi xi yi      − P mi xi zi
        Ib =  − P mi xi yi                   2   2
                                         i (xi + zi ) P−
                                        mP                  mi yi zi 
                     − mi xi zi        − mi yi zi       mi (x2i + yi2 )
                                    
                    Ixx Ixy Ixz
             =  Ixy Iyy Iyz  .
                    Ixz Iyz Izz

The summations can be replaced by volume integrals over the body B, using the
differential volume element dV , with the point masses mi replaced by a mass
density function ρ(x, y, z):
                                Z                            
                       Ixx =       (y 2 + z 2 )ρ(x, y, z) dV 
                                                             
                                                             
                                 B                           
                                                             
                                Z                            
                                                             
                                                             
                                                             
                       Iyy =       (x + z )ρ(x, y, z) dV 
                                      2     2
                                                             
                                                             
                                 B
                                                             
                                                             
                                Z                            
                                                             
                                                             
                                                             
                       Izz =       (x + y )ρ(x, y, z) dV 
                                      2     2                
                                                             
                                 B
\begin{equation}
Z
\label{eq:8.24}
\end{equation}
                                                             
                                                             
                       Ixy = − xyρ(x, y, z) dV               
                                                             
                                                             
                                                             
                                  ZB                         
                                                             
                                                             
                                                             
                       Ixz = − xzρ(x, y, z) dV               
                                                             
                                                             
                                                             
                                    B                        
                                                             
                                  Z                          
                                                             
                                                             
                                                             
                       Iyz = − yzρ(x, y, z) dV.              
                                              B

If the body has uniform density, Ib is determined exclusively by the shape of
the rigid body (see Figure 8.5).
    Given an inertia matrix Ib , the principal axes of inertia are given by
the eigenvectors and eigenvalues of Ib . Let v1 , v2 , v3 be the eigenvectors of
Ib and λ1 , λ2 , λ3 be the corresponding eigenvalues. Then the principal axes
of inertia are in the directions of v1 , v2 , v3 , and the scalar moments of inertia
about these axes, the principal moments of inertia, are λ1 , λ2 , λ3 > 0. One
principal axis maximizes the moment of inertia among all axes passing through
the center of mass, and another minimizes the moment of inertia. For bodies
with symmetry, often the principal axes of inertia are apparent. They may not
be unique; for a uniform-density solid sphere, for example, any three orthogonal
axes intersecting at the center of mass constitute a set of principal axes, and


286                                                          8.2. Dynamics of a Single Rigid Body

                    ẑ                                  ẑ                               ẑ

                                   h                                                c
                                           h       r
                              ŷ                                     ŷ       x̂     a        b     ŷ
            x̂                                     x̂
              w
  rectangular parallelepiped:               circular cylinder:                  ellipsoid:
        volume = abc,                       volume = πr2 h,                volume = 4πabc/3,
    Ixx = m(w2 + h2 )/12,                Ixx = m(3r2 + h2 )/12,            Ixx = m(b2 + c2 )/5,
     Iyy = m(`2 + h2 )/12,               Iyy = m(3r2 + h2 )/12,            Iyy = m(a2 + c2 )/5,
     Izz = m(`2 + w2 )/12                      Izz = mr2 /2                Izz = m(a2 + b2 )/5

\begin{figure}[h]
\centering
\includegraphics{figure_8_5.pdf}
\caption{The principal axes and the inertia about the principal axes for uniform-}
\label{fig:8_5}
\end{figure}
density bodies of mass m. Note that the x̂ and ŷ principal axes of the cylinder are not
unique.

the minimum principal moment of inertia is equal to the maximum principal
moment of inertia.
    If the principal axes of inertia are aligned with the axes of {b}, the off-
diagonal terms of Ib are all zero, and the eigenvalues are the scalar moments of
inertia Ixx , Iyy , and Izz about the x̂-, ŷ-, and ẑ-axes, respectively. In this case,
the equations of motion (8.23) simplify to
                                                              
                                Ixx ω̇x + (Izz − Iyy )ωy ωz
\begin{equation}
mb =  Iyy ω̇y + (Ixx − Izz )ωx ωz  ,
\label{eq:8.25}
\end{equation}
                                Izz ω̇z + (Iyy − Ixx )ωx ωy

where ωb = (ωx , ωy , ωz ). When possible, we choose the axes of {b} to be aligned
with the principal axes of inertia, in order to reduce the number of nonzero
entries in Ib and to simplify the equations of motion.
    Examples of common uniform-density solid bodies, their principal axes of
inertia, and the principal moments of inertia obtained by solving the inte-
grals (8.24), are given in Figure 8.5.
    An inertia matrix Ib can be expressed in a rotated frame {c} described by
the rotation matrix Rbc . Denoting this inertia matrix as Ic , and knowing that
the kinetic energy of the rotating body is independent of the chosen frame, we


Chapter 8. Dynamics of Open Chains                                                                  287

have
                            1 T                  1 T
                             ω Ic ωc       =       ω I b ωb
                            2 c                  2 b
                                                 1
                                           =       (Rbc ωc )T Ib (Rbc ωc )
                                                 2
                                                 1 T T
                                           =       ω (R Ib Rbc )ωc .
                                                 2 c bc
In other words,
                                                T
\begin{equation}
Ic = Rbc Ib Rbc .
\label{eq:8.26}
\end{equation}
If the axes of {b} are not aligned with the principal axes of inertia then we can
diagonalize the inertia matrix by expressing it instead in the rotated frame {c},
where the columns of Rbc correspond to the eigenvalues of Ib .
    Sometimes it is convenient to represent the inertia matrix in a frame at a
point not at the center of mass of the body, for example at a joint. Steiner’s
theorem can be stated as follows.
Theorem 8.2. The inertia matrix Iq about a frame aligned with {b}, but at a
point q = (qx , qy , qz ) in {b}, is related to the inertia matrix Ib calculated at the
center of mass by
\begin{equation}
Iq = Ib + m(q T qI − qq T ),
\label{eq:8.27}
\end{equation}
where I is the 3 × 3 identity matrix and m is the mass of the body.
    Steiner’s theorem is a more general statement of the parallel-axis theorem,
which states that the scalar inertia Id about an axis parallel to, but a distance
d from, an axis through the center of mass is related to the scalar inertia Icm
about the axis through the center of mass by
\begin{equation}
Id = Icm + md2 .
\label{eq:8.28}
\end{equation}
    Equations (8.26) and (8.27) are useful for calculating the inertia of a rigid
body consisting of component rigid bodies. First we calculate the inertia matri-
ces of the n component bodies in terms of frames at their individual centers of
mass. Then we choose a common frame {common} (e.g., at the center of mass
of the composite rigid body) and use Equations (8.26) and (8.27) to express each
inertia matrix in this common frame. Once the individual inertia matrices are
expressed in {common}, they can be summed to get the inertia matrix Icommon
for the composite rigid body.
    In the case of motion confined to the x̂–ŷ-plane, where ωb = (0, 0, ωz ) and
the inertia of the body about the ẑ-axis through the center of mass is given


288                                                      8.2. Dynamics of a Single Rigid Body

by the scalar Izz , the spatial rotational dynamics (8.23) reduces to the planar
rotational dynamics
                                   mz = Izz ω̇z ,
and the rotational kinetic energy is
                                                  1
                                           K=       Izz ωz2 .
                                                  2

\subsection{Twist–Wrench Formulation}
The linear dynamics (8.22) and the rotational dynamics (8.23) can be written
in the following combined form:
                                                       
     mb         Ib 0        ω̇b     [ωb ]  0       Ib 0        ωb
\begin{equation}
=                    +                                  ,
\label{eq:8.29}
\end{equation}
      fb         0 mI       v̇b       0   [ωb ]    0 mI        vb
where I is the 3 × 3 identity matrix. With the benefit of hindsight, and also
making use of the fact that [v]v = v × v = 0 and [v]T = −[v], we can write
Equation (8.29) in the following equivalent form:
                                                       
   mb         Ib 0          ω̇b      [ωb ] [vb ]     Ib 0      ωb
          =                      +
    fb         0 mI         v̇b       0    [ωb ]      0 mI     vb
                                            T              
              Ib 0          ω̇b      [ωb ]   0         Ib 0     ωb
\begin{equation}
=                      −                                    .
\label{eq:8.30}
\end{equation}
               0 mI         v̇b      [vb ] [ωb ]       0 mI     vb
    Written this way, each term can now be identified with six-dimensional spa-
tial quantities as follows:
  (a) The vectors (ωb , vb ) and (mb , fb ) can be respectively identified with the
      body twist Vb and body wrench Fb ,
                                                         
                                   ωb                   mb
\begin{equation}
Vb =           ,    Fb =          .
\label{eq:8.31}
\end{equation}
                                   vb                   fb

  (b) The spatial inertia matrix Gb ∈ R6×6 is defined as
                                             
                                      Ib 0
\begin{equation}
Gb =              .
\label{eq:8.32}
\end{equation}
                                      0 mI
       As an aside, the kinetic energy of the rigid body can be expressed in terms
       of the spatial inertia matrix as
                                          1 T       1         1
\begin{equation}
kinetic energy =        ω Ib ωb + mvbT vb = VbT Gb Vb .
\label{eq:8.33}
\end{equation}
                                          2 b       2         2


Chapter 8. Dynamics of Open Chains                                                                  289

  (c) The spatial momentum Pb ∈ R6 is defined as
                                             
                         Ib ωb     Ib 0        ωb
\begin{equation}
Pb =          =                   = Gb Vb .
\label{eq:8.34}
\end{equation}
                         mvb       0 mI        vb

   Observe that the term involving Pb in Equation (8.30) is left-multiplied by
the matrix                                  T
                                 [ωb ]  0
\begin{equation}
−                   .
\label{eq:8.35}
\end{equation}
                                 [vb ] [ωb ]
We now explain the origin and geometric significance of this matrix. First, recall
that the cross product of two vectors ω1 , ω2 ∈ R3 can be calculated, using the
skew-symmetric matrix notation, as follows:
\begin{equation}
[ω1 × ω2 ] = [ω1 ][ω2 ] − [ω2 ][ω1 ].
\label{eq:8.36}
\end{equation}
The matrix in (8.35) can be thought of as a generalization of the cross-product
operation to six-dimensional twists. Specifically, given two twists V1 = (ω1 , v1 )
and V2 = (ω2 , v2 ), we perform a calculation analogous to (8.36):
                                                                             
                            [ω1 ] v1        [ω2 ] v2          [ω2 ] v2     [ω1 ] v1
[V1 ][V2 ] − [V2 ][V1 ] =                                −
                             0      0          0    0           0     0     0    0
                                                                     
                            [ω1 ][ω2 ] − [ω2 ][ω1 ] [ω1 ]v2 − [ω2 ]v1
                        =
                                       0                    0
                           0        0
                                       
                            [ω ] v
                        =                ,
                             0     0
which can be written more compactly in vector form as
                       0                      
                        ω       [ω1 ]   0       ω2
                            =                        .
                        v0      [v1 ] [ω1 ]     v2
This generalization of the cross product to two twists V1 and V2 is called the
Lie bracket of V1 and V2 .
Definition 8.3. Given two twists V1 = (ω1 , v1 ) and V2 = (ω2 , v2 ), the Lie
bracket of V1 and V2 , written either as [adV1 ]V2 or adV1 (V2 ), is defined as
follows:                      
               [ω1 ]  0       ω2
\begin{equation}
= [adV1 ]V2 = adV1 (V2 ) ∈ R6 ,
\label{eq:8.37}
\end{equation}
               [v1 ] [ω1 ]    v2
where                                                      
                                                [ω]    0
\begin{equation}
[adV ] =                       ∈ R6×6 .
\label{eq:8.38}
\end{equation}
                                                [v]   [ω]


290                                                      8.2. Dynamics of a Single Rigid Body

Definition 8.4. Given a twist V = (ω, v) and a wrench F = (m, f ), define the
mapping
                                  T                        
     T             T      [ω] 0          m        −[ω]m − [v]f
\begin{equation}
adV (F) = [adV ] F =                       =                   .
\label{eq:8.39}
\end{equation}
                          [v] [ω]        f           −[ω]f
    Using the notation and definitions above, the dynamic equations for a single
rigid body can now be written as

                                 Fb    = Gb V̇b − adT
                                                    Vb (Pb )
\begin{equation}
= Gb V̇b − [adVb ]T Gb Vb .
\label{eq:8.40}
\end{equation}
Note the analogy between Equation (8.40) and the moment equation for a ro-
tating rigid body:
\begin{equation}
mb = Ib ω̇b − [ωb ]T Ib ωb .
\label{eq:8.41}
\end{equation}
Equation (8.41) is simply the rotational component of (8.40).

\subsection{Dynamics in Other Frames}
The derivation of the dynamic equations (8.40) relies on the use of a center-of-
mass frame {b}. It is straightforward to express the dynamics in other frames,
however. Let’s call one such frame {a}.
    Since the kinetic energy of the rigid body must be independent of the frame
of representation,
                          1 T       1
                           V Ga Va = VbT Gb Vb
                          2 a       2
                                    1
                                   = ([AdTba ]Va )T Gb [AdTba ]Va
                                    2
                                    1
                                   = VaT [AdTba ]T Gb [AdTba ] Va ;
                                    2    |        {z        }
                                                            Ga

for the adjoint representation Ad (see Definition 3.20). In other words, the
spatial inertia matrix Ga in {a} is related to Gb by
\begin{equation}
Ga = [AdTba ]T Gb [AdTba ].
\label{eq:8.42}
\end{equation}
This is a generalization of Steiner’s theorem.
   Using the spatial inertia matrix Ga , the equations of motion (8.40) in the
{b} frame can be expressed equivalently in the {a} frame as
\begin{equation}
Fa = Ga V̇a − [adVa ]T Ga Va ,
\label{eq:8.43}
\end{equation}
Chapter 8. Dynamics of Open Chains                                                                  291

where Fa and Va are the wrench and twist written in the {a} frame. (See
Exercise 8.3.) Thus the form of the equations of motion is independent of the
frame of representation.

\section{Newton–Euler Inverse Dynamics}
We now consider the inverse dynamics problem for an n-link open chain con-
nected by one-dof joints. Given the joint positions θ ∈ Rn , velocities θ̇ ∈ Rn ,
and accelerations θ̈ ∈ Rn , the objective is to calculate the right-hand side of the
dynamics equation
                               τ = M (θ)θ̈ + h(θ, θ̇).
The main result is a recursive inverse dynamics algorithm consisting of a forward
and a backward iteration stage. In the former, the positions, velocities, and
accelerations of each link are propagated from the base to the tip while in
the backward iterations the forces and moments experienced by each link are
propagated from the tip to the base.

\subsection{Derivation}
A body-fixed reference frame {i} is attached to the center of mass of each link
i, i = 1, . . . , n. The base frame is denoted {0}, and a frame at the end-effector
is denoted {n + 1}. This frame is fixed in {n}.
     When the manipulator is at the home position, with all joint variables zero,
we denote the configuration of frame {j} in {i} as Mi,j ∈ SE(3), and the
configuration of {i} in the base frame {0} using the shorthand Mi = M0,i .
With these definitions, Mi−1,i and Mi,i−1 can be calculated as
                            −1
                  Mi−1,i = Mi−1 Mi              and        Mi,i−1 = Mi−1 Mi−1 .

    The screw axis for joint i, expressed in the link frame {i}, is Ai . This same
screw axis is expressed in the space frame {0} as Si , where the two are related
by
                                 Ai = AdM −1 (Si ).
                                                      i

Defining Ti,j ∈ SE(3) to be the configuration of frame {j} in {i} for arbitrary
joint variables θ then Ti−1,i (θi ), the configuration of {i} relative to {i − 1} given
                                            −1
the joint variable θi , and Ti,i−1 (θi ) = Ti−1,i (θi ) are calculated as

         Ti−1,i (θi ) = Mi−1,i e[Ai ]θi         and        Ti,i−1 (θi ) = e−[Ai ]θi Mi,i−1 .

    We further adopt the following notation:


292                                                       8.3. Newton–Euler Inverse Dynamics

  (a) The twist of link frame {i}, expressed in frame-{i} coordinates, is denoted
      Vi = (ωi , vi ).
  (b) The wrench transmitted through joint i to link frame {i}, expressed in
      frame-{i} coordinates, is denoted Fi = (mi , fi ).
  (c) Let Gi ∈ R6×6 denote the spatial inertia matrix of link i, expressed relative
      to link frame {i}. Since we are assuming that all link frames are situated
      at the link center of mass, Gi has the block-diagonal form
                                                  
                                          Ii   0
\begin{equation}
Gi =               ,
\label{eq:8.44}
\end{equation}
                                          0 mi I
       where Ii denotes the 3 × 3 rotational inertia matrix of link i and mi is the
       link mass.
    With these definitions, we can recursively calculate the twist and acceleration
of each link, moving from the base to the tip. The twist Vi of link i is the sum
of the twist of link i − 1, but expressed in {i}, and the added twist due to the
joint rate θ̇i :
\begin{equation}
Vi = Ai θ˙i + [AdTi,i−1 ]Vi−1 .
\label{eq:8.45}
\end{equation}
    The accelerations V̇i can also be found recursively. Taking the time derivative
of Equation (8.45), we get
                                              d             
                                                 [AdTi,i−1 ] Vi−1 .
\begin{equation}
V̇i = Ai θ̈i + [AdTi,i−1 ]V̇i−1 +
\label{eq:8.46}
\end{equation}
                                              dt
To calculate the final term in this equation, we express Ti,i−1 and Ai as
                                                             
                           Ri,i−1 p                           ω
               Ti,i−1 =                     and     Ai =          .
                             0      1                         v
Then
                                                       
  d                     d          Ri,i−1         0
     [AdTi,i−1 ] Vi−1 =                                     Vi−1
  dt                    dt       [p]Ri,i−1 Ri,i−1
                                                                              
                                      −[ω θ̇i ]Ri,i−1                 0
                      =                                                          Vi−1
                           −[v θ̇i ]Ri,i−1 − [ω θ̇i ][p]Ri,i−1 −[ω θ̇i ]Ri,i−1
                                                                      
                           −[ω θ̇i ]      0            Ri,i−1     0
                      =                                                    Vi−1
                           −[v θ̇i ] −[ω θ̇i ]        [p]Ri,i−1 Ri,i−1
                        |            {z          }|           {z         }
                                        −[adA θ̇ ]                  [AdTi,i−1 ]
                                                 i i

                             = −[adAi θ̇i ]Vi
                             = [adVi ]Ai θ̇i .


Chapter 8. Dynamics of Open Chains                                                                  293

                    joint
                    axis i              V̇i
                                                              Vi

                                                 {i}

              Fi
                                                                    − AdTTi+1,i (Fi+1 )

\begin{figure}[h]
\centering
\includegraphics{figure_8_6.pdf}
\caption{Free-body diagram illustrating the moments and forces exerted on link}
\label{fig:8_6}
\end{figure}
i.

Substituting this result into Equation (8.46), we get
\begin{equation}
V̇i = Ai θ̈i + [AdTi,i−1 ]V̇i−1 + [adVi ]Ai θ̇i ,
\label{eq:8.47}
\end{equation}
i.e., the acceleration of link i is the sum of three components: a component due
to the joint acceleration θ̈i , a component due to the acceleration of link i − 1
expressed in {i}, and a velocity-product component.
     Once we have determined all the link twists and accelerations moving out-
ward from the base, we can calculate the joint torques or forces by moving
inward from the tip. The rigid-body dynamics (8.40) tells us the total wrench
that acts on link i given Vi and V̇i . Furthermore, the total wrench acting on
link i is the sum of the wrench Fi transmitted through joint i and the wrench
applied to the link through joint i + 1 (or, for link n, the wrench applied to the
link by the environment at the end-effector frame {n + 1}), expressed in the
frame i. Therefore, we have the equality

                         Gi V̇i − adT                    T
\begin{equation}
Vi (Gi Vi ) = Fi − AdTi+1,i (Fi+1 );
\label{eq:8.48}
\end{equation}
see Figure 8.6. Solving from the tip toward the base, at each link i we solve
for the only unknown in Equation (8.48): Fi . Since joint i has only one degree
of freedom, five dimensions of the six-vector Fi are provided “for free” by the
structure of the joint, and the actuator only has to provide the scalar force or
torque in the direction of the joint’s screw axis:
\begin{equation}
τi = FiT Ai .
\label{eq:8.49}
\end{equation}
Equation (8.49) provides the torques required at each joint, solving the inverse
dynamics problem.


294                                                    8.4. Dynamic Equations in Closed Form

\subsection{Newton-Euler Inverse Dynamics Algorithm}
Initialization Attach a frame {0} to the base, frames {1} to {n} to the centers
of mass of links {1} to {n}, and a frame {n + 1} at the end-effector, fixed in the
frame {n}. Define Mi,i−1 to be the configuration of {i − 1} in {i} when θi = 0.
Let Ai be the screw axis of joint i expressed in {i}, and Gi be the 6 × 6 spatial
inertia matrix of link i. Define V0 to be the twist of the base frame {0} expressed
in {0} coordinates. (This quantity is typically zero.) Let g ∈ R3 be the gravity
vector expressed in base-frame coordinates, and define V̇0 = (ω̇0 , v̇0 ) = (0, −g).
(Gravity is treated as an acceleration of the base in the opposite direction.)
Define Fn+1 = Ftip = (mtip , ftip ) to be the wrench applied to the environment
by the end-effector, expressed in the end-effector frame {n + 1}.

Forward iterations Given θ, θ̇, θ̈, for i = 1 to n do
\begin{equation}
Ti,i−1      = e−[Ai ]θi Mi,i−1 ,
\label{eq:8.50}
\end{equation}
\begin{equation}
Vi     =     AdTi,i−1 (Vi−1 ) + Ai θ̇i ,
\label{eq:8.51}
\end{equation}
\begin{equation}
V̇i    =     AdTi,i−1 (V̇i−1 ) + adVi (Ai )θ̇i + Ai θ̈i .
\label{eq:8.52}
\end{equation}
Backward iterations For i = n to 1 do

                      Fi       =    AdT                           T
\begin{equation}
Ti+1,i (Fi+1 ) + Gi V̇i − adVi (Gi Vi ),
\label{eq:8.53}
\end{equation}
\begin{equation}
τi      = FiT Ai .
\label{eq:8.54}
\end{equation}
\section{Dynamic Equations in Closed Form}
In this section we show how the equations in the recursive inverse dynamics
algorithm can be organized into a closed-form set of dynamics equations τ =
M (θ)θ̈ + c(θ, θ̇) + g(θ).
   Before doing so, we prove our earlier assertion that the total kinetic energy
K of the robot can be expressed as K = 21 θ̇T M (θ)θ̇. We do so by noting that K
can be expressed as the sum of the kinetic energies of each link:
                                                   n
                                               1X T
\begin{equation}
K=         V Gi Vi ,
\label{eq:8.55}
\end{equation}
                                               2 i=1 i

where Vi is the twist of link frame {i} and Gi is the spatial inertia matrix
of link i as defined by Equation (8.32) (both are expressed in link-frame-{i}
coordinates). Let T0i (θ1 , . . . , θi ) denote the forward kinematics from the base


Chapter 8. Dynamics of Open Chains                                                                  295

frame {0} to link frame {i}, and let Jib (θ) denote the body Jacobian obtained
        −1
from T0i   Ṫ0i . Note that Jib as defined is a 6 × i matrix; we turn it into a 6 × n
matrix by filling in all entries of the last n − i columns with zeros. With this
definition of Jib , we can write

                                Vi = Jib (θ)θ̇,          i = 1, . . . , n.

The kinetic energy can then be written
                                             n
                                                                      !
                                 1           X
                              K = θ̇T                T
\begin{equation}
Jib (θ)Gi Jib (θ)     θ̇.
\label{eq:8.56}
\end{equation}
                                 2            i=1

The term inside the parentheses is precisely the mass matrix M (θ):
                                              n
                                              X
                                                     T
\begin{equation}
M (θ) =           Jib (θ)Gi Jib (θ).
\label{eq:8.57}
\end{equation}
                                              i=1

   We now return to the original task of deriving a closed-form set of dynamic
equations. We start by defining the following stacked vectors:
                                         
                                       V1
                                         
\begin{equation}
V =  ...  ∈ R6n ,
\label{eq:8.58}
\end{equation}
                                                    Vn
                                                 
                                              F1
                                                 
\begin{equation}
F     =  ...  ∈ R6n .
\label{eq:8.59}
\end{equation}
                                                    Fn


296                                                   8.4. Dynamic Equations in Closed Form

Further, define the following matrices:
                                                      
                       A1 0 · · · 0
                    0 A2 · · · 0                      
                                                      
\begin{equation}
A =  .           .. ..     ..               ∈ R6n×n ,
\label{eq:8.60}
\end{equation}
                    ..       .    .    .              
                              0    ···    ···    An
                                                     
                       G1          0     ···     0
                      0           G2    ···     0    
                                                     
\begin{equation}
G    =  .            ..   ..      ..    ∈ R6n×6n ,
\label{eq:8.61}
\end{equation}
                      ..            .      .     .   
                             0     ···   ···    Gn
                                                             
                       [adV1 ]    0         ···            0
                        0     [adV2 ] · · ·               0  
                                                             
\begin{equation}
[adV ]   =     ..       ..       ..             ..  ∈ R6n×6n ,
\label{eq:8.62}
\end{equation}
                          .        .           .           . 
                     0           ···        · · · [adVn ]
                                                                       
                   [adA1 θ̇1 ]           0           ···       0
                    0          [adA2 θ̇2 ] · · ·            0        
                                                                       
         adAθ̇ =       ..               ..          . .        .
\begin{equation}
.        ∈ R6n×6n ,
\label{eq:8.63}
\end{equation}
                        .                .              .      .       
                       0               ···           · · · [adAn θ̇n ]
                                                                            
                      0               0           ···         0           0
                  [AdT21 ]           0           ···         0           0 
                                                                            
                     0          [AdT32 ] · · ·               0           0 
         W(θ) =                                                              ∈ R6n×6n .(8.64)
                     ..              ..          . ..        .
                                                              ..          .. 
                                                                          .
                      .               .                                     
                                                                     
                      0               0           · · · AdTn,n−1          0

We write W(θ) to emphasize the dependence of W on θ. Finally, define the


Chapter 8. Dynamics of Open Chains                                                                  297

following stacked vectors:
                                                  
                                       AdT10 (V0 )
                                          0       
                                                  
\begin{equation}
Vbase    =      ..       ∈ R6n ,
\label{eq:8.65}
\end{equation}
                                           .      
                                                  0
                                                   
                                       AdT10 (V̇0 )
                                          0        
                                                   
\begin{equation}
V̇base   =      ..        ∈ R6n ,
\label{eq:8.66}
\end{equation}
                                           .       
                                                  0
                                                                  
                                                      0
                                                     ..           
                                                      .           
\begin{equation}
Ftip    =                               ∈ R6n .
\label{eq:8.67}
\end{equation}
                                                     0            
                                             AdT
                                               Tn+1,n (Fn+1 )

Note that A ∈ R6n×n and G ∈ R6n×6n are constant block-diagonal matrices,
in which A contains only the kinematic parameters while G contains only the
mass and inertial parameters for each link.
   With the above definitions, our earlier recursive inverse dynamics algorithm
can be assembled into the following set of matrix equations:
\begin{equation}
V    = W(θ)V + Aθ̇ + Vbase ,
\label{eq:8.68}
\end{equation}
\begin{equation}
V̇   = W(θ)V̇ + Aθ̈ − [adAθ̇ ](W(θ)V + Vbase ) + V̇base ,
\label{eq:8.69}
\end{equation}
\begin{equation}
F     = W T (θ)F + G V̇ − [adV ]T GV + Ftip ,
\label{eq:8.70}
\end{equation}
\begin{equation}
τ     = AT F.
\label{eq:8.71}
\end{equation}
The matrix W(θ) has the property that W n (θ) = 0 (such a matrix is said to be
nilpotent of order n), and one consequence verifiable through direct calculation
is that (I −W(θ))−1 = I +W(θ)+· · ·+W n−1 (θ). Defining L(θ) = (I −W(θ))−1 ,
it can further be verified via direct calculation that
                                                         
                     I          0         0      ··· 0
                [AdT21 ]       I         0      ··· 0 
                                                         
                [AdT31 ] [AdT32 ]        I      ··· 0 
\begin{equation}
L(θ) =                                            ∈ R6n×6n .
\label{eq:8.72}
\end{equation}
                     ..        ..         ..    . .   .. 
                      .         .          .        . . 
                       [AdTn1 ]     [AdTn2 ]    [AdTn3 ] · · ·      I

We write L(θ) to emphasize the dependence of L on θ. The earlier matrix


298                                                    8.5. Forward Dynamics of Open Chains

equations can now be reorganized as follows:
                                 
\begin{equation}
V = L(θ) Aθ̇ + Vbase ,
\label{eq:8.73}
\end{equation}
                                                                
\begin{equation}
V̇ = L(θ) Aθ̈ + [adAθ̇ ]W(θ)V + [adAθ̇ ]Vbase + V̇base ,
\label{eq:8.74}
\end{equation}
                                            
\begin{equation}
F = LT (θ) G V̇ − [adV ]T GV + Ftip ,
\label{eq:8.75}
\end{equation}
                           T
\begin{equation}
τ    = A F.
\label{eq:8.76}
\end{equation}
If the robot applies an external wrench Ftip at the end-effector, this can be
included into the dynamics equation
\begin{equation}
τ = M (θ)θ̈ + c(θ, θ̇) + g(θ) + J T (θ)Ftip ,
\label{eq:8.77}
\end{equation}
where J(θ) denotes the Jacobian of the forward kinematics expressed in the
same reference frame as Ftip , and
\begin{equation}
M (θ) = AT LT (θ)GL(θ)A,
\label{eq:8.78}
\end{equation}
                                                              
\begin{equation}
c(θ, θ̇) = −AT LT (θ) GL(θ) [adAθ̇ ]W(θ) + [adV ]T G L(θ)Aθ̇,
\label{eq:8.79}
\end{equation}
\begin{equation}
g(θ)     = AT LT (θ)GL(θ)V̇base .
\label{eq:8.80}
\end{equation}
\section{Forward Dynamics of Open Chains}
The forward dynamics problem involves solving
\begin{equation}
M (θ)θ̈ = τ (t) − h(θ, θ̇) − J T (θ)Ftip
\label{eq:8.81}
\end{equation}
for θ̈, given θ, θ̇, τ , and the wrench Ftip applied by the end-effector (if ap-
plicable). The term h(θ, θ̇) can be computed by calling the inverse dynamics
algorithm with θ̈ = 0 and Ftip = 0. The inertia matrix M (θ) can be computed
using Equation (8.57). An alternative is to use n calls of the inverse dynamics
algorithm to build M (θ) column by column. In each of the n calls, set g = 0,
θ̇ = 0, and Ftip = 0. In the first call, the column vector θ̈ is all zeros except
for a 1 in the first row. In the second call, θ̈ is all zeros except for a 1 in the
second row, and so on. The τ vector returned by the ith call is the ith column
of M (θ), and after n calls the n × n matrix M (θ) is constructed.
    With M (θ), h(θ, θ̇), and Ftip we can use any efficient algorithm for solving
Equation (8.81), which is of the form M θ̈ = b, for θ̈.
    The forward dynamics can be used to simulate the motion of the robot given
its initial state, the joint forces–torques τ (t), and an optional external wrench


Chapter 8. Dynamics of Open Chains                                                                  299

Ftip (t), for t ∈ [0, tf ]. First define the function ForwardDynamics returning the
solution to Equation (8.81), i.e.,

                               θ̈ = ForwardDynamics(θ, θ̇, τ, Ftip ).

Defining the variables q1 = θ, q2 = θ̇, the second-order dynamics (8.81) can be
converted to two first-order differential equations,

                         q̇1    = q2 ,
                         q̇2    = ForwardDynamics(q1 , q2 , τ, Ftip ).

The simplest method for numerically integrating a system of first-order differ-
ential equations of the form q̇ = f (q, t), q ∈ Rn , is the first-order Euler iteration

                                  q(t + δt) = q(t) + δtf (q(t), t),

where the positive scalar δt denotes the timestep. The Euler integration of the
robot dynamics is thus

              q1 (t + δt) = q1 (t) + q2 (t)δt,
              q2 (t + δt) = q2 (t) + ForwardDynamics(q1 , q2 , τ, Ftip )δt.

Given a set of initial values for q1 (0) = θ(0) and q2 (0) = θ̇(0), the above equa-
tions can be iterated forward in time to obtain the motion θ(t) = q1 (t) numeri-
cally.

Euler Integration Algorithm for Forward Dynamics
    • Inputs: The initial conditions θ(0) and θ̇(0), the input torques τ (t) and
      wrenches at the end-effector Ftip (t) for t ∈ [0, tf ], and the number of
      integration steps N .
    • Initialization: Set the timestep δt = tf /N , and set θ[0] = θ(0), θ̇[0] =
      θ̇(0).
    • Iteration: For k = 0 to N − 1 do

                     θ̈[k]     = ForwardDynamics(θ[k], θ̇[k], τ (kδt), Ftip (kδt)),
                θ[k + 1] = θ[k] + θ̇[k]δt,
                θ̇[k + 1] = θ̇[k] + θ̈[k]δt.

    • Output: The joint trajectory θ(kδt) = θ[k], θ̇(kδt) = θ̇[k], k = 0, . . . , N .


300                                                            8.6. Dynamics in the Task Space

    The result of the numerical integration converges to the theoretical result
as the number of integration steps N goes to infinity. Higher-order numerical
integration schemes, such as fourth-order Runge–Kutta, can yield a closer ap-
proximation with fewer computations than the simple first-order Euler method.

\section{Dynamics in the Task Space}
In this section we consider how the dynamic equations change under a trans-
formation to coordinates of the end-effector frame (task-space coordinates). To
keep things simple we consider a six-degree-of-freedom open chain with joint-
space dynamics
\begin{equation}
τ = M (θ)θ̈ + h(θ, θ̇),           θ ∈ R6 , τ ∈ R6 .
\label{eq:8.82}
\end{equation}
We also ignore, for the time being, any end-effector forces Ftip . The twist
V = (ω, v) of the end-effector is related to the joint velocity θ̇ by
\begin{equation}
V = J(θ)θ̇,
\label{eq:8.83}
\end{equation}
with the understanding that V and J(θ) are always expressed in terms of the
same reference frame. The time derivative V̇ is then
                                              ˙ θ̇ + J(θ)θ̈.
\begin{equation}
V̇ = J(θ)
\label{eq:8.84}
\end{equation}
At configurations θ where J(θ) is invertible, we have
\begin{equation}
θ̇   = J −1 V,
\label{eq:8.85}
\end{equation}
                                  θ̈                     ˙ −1 V.
\begin{equation}
= J −1 V̇ − J −1 JJ
\label{eq:8.86}
\end{equation}
Substituting for θ̇ and θ̈ in Equation (8.82) leads to
                                                
                                           ˙ −1 V + h(θ, J −1 V).
\begin{equation}
τ = M (θ) J −1 V̇ − J −1 JJ
\label{eq:8.87}
\end{equation}
Let J −T denote (J −1 )T = (J T )−1 . Pre-multiply both sides by J −T to get

                      J −T τ                                     ˙ −1 V
                                 = J −T M J −1 V̇ − J −T M J −1 JJ
\begin{equation}
−T         −1
\label{eq:8.88}
\end{equation}
                                   + J h(θ, J V).

Expressing J −T τ as the wrench F, the above can be written
\begin{equation}
F = Λ(θ)V̇ + η(θ, V),
\label{eq:8.89}
\end{equation}
Chapter 8. Dynamics of Open Chains                                                                  301

where
\begin{equation}
Λ(θ) = J −T M (θ)J −1 ,
\label{eq:8.90}
\end{equation}
                                                          ˙ −1 V.
\begin{equation}
η(θ, V) = J −T h(θ, J −1 V) − Λ(θ)JJ
\label{eq:8.91}
\end{equation}
These are the dynamic equations expressed in end-effector frame coordinates.
If an external wrench F is applied to the end-effector frame then, assuming the
actuators provide zero forces and torques, the motion of the end-effector frame
is governed by these equations.
    Note that J(θ) must be invertible (i.e., there must be a one-to-one mapping
between joint velocities and end-effector twists) in order to derive the task-
space dynamics above. Also note the dependence of Λ(θ) and η(θ, V) on θ. In
general, we cannot replace the dependence on θ by a dependence on the end-
effector configuration X because there may be multiple solutions to the inverse
kinematics, and the dynamics depends on the specific joint configuration θ.

\section{Constrained Dynamics}
Now consider the case where the n-joint robot is subject to a set of k holonomic
or nonholonomic Pfaffian velocity constraints of the form
\begin{equation}
A(θ)θ̇ = 0,         A(θ) ∈ Rk×n .
\label{eq:8.92}
\end{equation}
(See Section 2.4 for an introduction to Pfaffian constraints.) Such constraints
can come from loop-closure constraints; for example, the motion of an end-
effector rigidly holding a door handle is subject to k = 5 constraints due to the
hinges of the door. As another example, a robot writing with a pen is subject
to a single constraint that keeps the height of the tip of the pen above the paper
at zero. In any case, we assume that the constraints do no work on the robot,
i.e., the generalized forces τcon due to the constraints satisfy
                                              T
                                             τcon θ̇ = 0.

This assumption means that τcon must be a linear combination of the columns
of AT (θ), i.e., τcon = AT (θ)λ for some λ ∈ Rk , since these are the generalized
forces that do no work when θ̇ is subject to the constraints (8.92):

                     (AT (θ)λ)T θ̇ = λT A(θ)θ̇ = 0              for all λ ∈ Rk .

For the writing-robot example, the assumption that the constraint is workless
means that there can be no friction between the pen and the paper.


302                                                                   8.7. Constrained Dynamics

   Adding the constraint forces AT (θ)λ to the equations of motion, we can
write the n + k constrained equations of motion in the n + k unknowns {θ̈, λ}
(the forward dynamics) or in the n + k unknowns {τ, λ} (the inverse dynamics):
\begin{equation}
τ      = M (θ)θ̈ + h(θ, θ̇) + AT (θ)λ,
\label{eq:8.93}
\end{equation}
\begin{equation}
A(θ)θ̇     = 0,
\label{eq:8.94}
\end{equation}
where λ is a set of Lagrange multipliers and AT (θ)λ are the joint forces and
torques that act against the constraints. From these equations, it should be
clear that the robot has n − k velocity freedoms and k “force freedoms” – the
robot is free to create any generalized force of the form AT (θ)λ. (For the writing
robot, there is an equality constraint: the robot can only apply pushing forces
into the paper and table, not pulling forces.)
    Often it is convenient to reduce the n + k equations in n + k unknowns to n
equations in n unknowns, without explicitly calculating the Lagrange multipliers
λ. To do this, we can solve for λ in terms of the other quantities and substitute
our solution into Equation (8.93). Since the constraints are satisfied at all times,
the time rate of change of the constraints satisfies
\begin{equation}
Ȧ(θ)θ̇ + A(θ)θ̈ = 0.
\label{eq:8.95}
\end{equation}
Assuming that M (θ) and A(θ) are full rank, we can solve Equation (8.93) for
θ̈, substitute into Equation (8.95), and omit the dependences on θ and θ̇ for
conciseness, to get
\begin{equation}
Ȧθ̇ + AM −1 (τ − h − AT λ) = 0.
\label{eq:8.96}
\end{equation}
Using Aθ̈ = −Ȧθ̇, after some manipulation, we obtain
\begin{equation}
λ = (AM −1 AT )−1 (AM −1 (τ − h) − Aθ̈).
\label{eq:8.97}
\end{equation}
Combining Equation (8.97) with Equation (8.93) and manipulating further, we
get
\begin{equation}
P τ = P (M θ̈ + h)
\label{eq:8.98}
\end{equation}
where
\begin{equation}
P = I − AT (AM −1 AT )−1 AM −1
\label{eq:8.99}
\end{equation}
and where I is the n × n identity matrix. The n × n matrix P (θ) is of rank
n − k, and maps the joint generalized forces τ to P (θ)τ , projecting away the
generalized force components that act on the constraints while retaining the
generalized forces that do work on the robot. The complementary projection,
I − P (θ), maps τ to (I − P (θ))τ , the joint forces that act on the constraints and
do no work on the robot.


Chapter 8. Dynamics of Open Chains                                                                  303

    Using P (θ), we can solve the inverse dynamics by evaluating the right-hand
side of Equation (8.98). The solution P (θ)τ is a set of generalized joint forces
that achieves the feasible components of the desired joint acceleration θ̈. To
this solution we can add any constraint forces AT (θ)λ without changing the
acceleration of the robot.
    We can rearrange Equation (8.98) into the form
\begin{equation}
Pθ̈ θ̈ = Pθ̈ M −1 (τ − h),
\label{eq:8.100}
\end{equation}
where
\begin{equation}
Pθ̈ = M −1 P M = I − M −1 AT (AM −1 AT )−1 A.
\label{eq:8.101}
\end{equation}
The rank-(n−k) projection matrix Pθ̈ (θ) ∈ Rn×n projects away the components
of an arbitrary joint acceleration θ̈ that violate the constraints, leaving only the
components Pθ̈ (θ)θ̈ that satisfy the constraints. To solve the forward dynamics,
we evaluate the right-hand side of Equation (8.100), and the resulting Pθ̈ θ̈ is
the set of joint accelerations.
    In Section 11.6 we discuss the related topic of hybrid motion–force control, in
which the goal at each instant is to simultaneously achieve a desired acceleration
satisfying Aθ̈ = 0 (consisting of n − k motion freedoms) and a desired force
against the k constraints. In that section we use the task-space dynamics to
represent the task-space end-effector motions and forces more naturally.

\section{Robot Dynamics in the URDF}
As described in Section 4.2 and illustrated in the UR5 Universal Robot Descrip-
tion Format file, the inertial properties of link i are described in the URDF by
the link elements mass, origin (the position and orientation of the center-of-
mass frame relative to a frame attached at joint i), and inertia, which specifies
the six elements of the symmetric rotational inertia matrix on or above the di-
agonal. To fully write the robot’s dynamics, for joint i we need in addition the
joint element origin, specifying the position and orientation of link i’s joint
frame relative to link (i − 1)’s joint frame when θi = 0, and the element axis,
which specifies the axis of motion of joint i. We leave to the exercises the trans-
lation of these elements into the quantities needed for the Newton–Euler inverse
dynamics algorithm.

\section{Actuation, Gearing, and Friction}
Until now we have been assuming the existence of actuators that directly provide
commanded forces and torques. In practice there are many types of actuators


304                                                      8.9. Actuation, Gearing, and Friction

(e.g., electric, hydraulic, and pneumatic) and mechanical power transformers
(e.g., gearheads), and the actuators can be located at the joints themselves or
remotely, with mechanical power transmitted by cables or timing belts. Each
combination of these has its own characteristics that can play a significant role
in the “extended dynamics” mapping the actual control inputs (e.g., the current
requested of amplifiers connected to electric motors) to the motion of the robot.
    In this section we provide an introduction to some of the issues associated
with one particular, and common, configuration: geared DC electric motors
located at each joint. This is the configuration used in the Universal Robots
UR5, for example.
    Figure 8.7 shows the electrical block diagram for a typical n-joint robot
driven by DC electric motors. For concreteness, we assume that each joint is
revolute. A power supply converts the wall AC voltage to a DC voltage to power
the amplifier associated with each motor. A control box takes user input, for
example in the form of a desired trajectory, as well as position feedback from
encoders located at each joint. Using the desired trajectory, a model of the
robot’s dynamics, and the measured error in the current robot state relative
to the desired robot state, the controller calculates the torque required of each
actuator. Since DC electric motors nominally provide a torque proportional to
the current through the motor, this torque command is equivalent to a current
command. Each motor amplifier then uses a current sensor (shown as external to
the amplifier in Figure 8.7, but in reality internal to the amplifier) to continually
adjust the voltage across the motor to try to achieve the requested current.3
The motion of the motor is sensed by the motor encoder, and the position
information is sent back to the controller.
    The commanded torque is typically updated at around 1000 times per second
(1 kHz), and the amplifier’s voltage control loop may be updated at a rate ten
times that or more.
    Figure 8.8 is a conceptual representation of the motor and other components
for a single axis. The motor has a single shaft extending from both ends of
the motor: one end drives a rotary encoder, which measures the position of
the joint, and the other end becomes the input to a gearhead. The gearhead
increases the torque while reducing the speed, since most DC electric motors
with an appropriate power rating provide torques that are too low to be useful
for robotics applications. The purpose of the bearing is to support the gearhead
output, freely transmitting torques about the gearhead axis while isolating the
gearhead (and motor) from wrench components due to link i + 1 in the other
five directions. The outer cases of the encoder, motor, gearhead, and bearing
   3 The voltage is typically a time-averaged voltage achieved by the duty cycle of a voltage

rapidly switching between a maximum positive voltage and a maximum negative voltage.


Chapter 8. Dynamics of Open Chains                                                                  305

                                             position feedback

                                                              I1                 I1
                    torque/                                         current            motor 1
                                                 amp 1
                     current                                         sensor           + encoder
  controller,      commands
  dynamics
                                                              I2                 I2
    model                                                                              motor 2
                                                                    current
                                                 amp 2
                                                                     sensor           + encoder
  user input
                                                                      .
                       DC                                             .
    power            voltage                                          .
    supply
                                                              In                In
                                                                    current            motor n
                                                 amp n
      AC                                                             sensor           + encoder
    voltage

\begin{figure}[h]
\centering
\includegraphics{figure_8_7.pdf}
\caption{A block diagram of a typical n-joint robot. The bold lines correspond}
\label{fig:8_7}
\end{figure}
to high-power signals while the thin lines correspond to communication signals.

                                current

   position
  feedback

                encoder               motor                   gearhead        bearing      link i+1

\begin{figure}[h]
\centering
\includegraphics{figure_8_8.pdf}
\caption{The outer cases of the encoder, motor, gearhead, and bearing are fixed}
\label{fig:8_8}
\end{figure}
in link i, while the gearhead output shaft supported by the bearing is fixed in link
i + 1.

are all fixed relative to each other and to link i. It is also common for the motor
to have some kind of brake, not shown.

\subsection{DC Motors and Gearing}
A DC motor consists of a stator and a rotor that rotates relative to the stator.
DC electric motors create torque by sending current through windings in a


306                                                      8.9. Actuation, Gearing, and Friction

magnetic field created by permanent magnets, where the magnets are attached
to the stator and the windings are attached to the rotor, or vice versa. A DC
motor has multiple windings, some of which are energized and some of which
are inactive at any given time. The windings that are energized are chosen as a
function of the angle of the rotor relative to the stator. This “commutation” of
the windings occurs mechanically using brushes (brushed motors) or electrically
using control circuitry (brushless motors). Brushless motors have the advantage
of no brush wear and higher continuous torque, since the windings are typically
attached to the motor housing where the heat due to the resistance of the
windings can be more easily dissipated. In our basic introduction to DC motor
modeling, we do not distinguish between brushed and brushless motors.
    Figure 8.9 shows a brushed DC motor with an encoder and a gearhead.
    The torque τ , measured in newton-meters (N m), created by a DC motor is
governed by the equation
                                     τ = kt I,
where I, measured in amps (A), is the current through the windings. The
constant kt , measured in newton-meters per amp (N m/A), is called the torque
constant. The power dissipated as heat by the windings, measured in watts
(W), is governed by
                                 Pheat = I 2 R,
where R is the resistance of the windings in ohms (Ω). To keep the motor wind-
ings from overheating, the continuous current flowing through the motor must
be limited. Accordingly, in continuous operation, the motor torque must be kept
below a continuous-torque limit τcont determined by the thermal properties of
the motor.
    A simplified model of a DC motor, where all units are in the SI system, can
be derived by equating the electrical power consumed by the motor Pelec = IV
in watts (W) to the mechanical power Pmech = τ w (also in W) and other power
produced by the motor,
                                       dI
          IV = τ w + I 2 R + LI           + friction and other power-loss terms,
                                       dt
where V is the voltage applied to the motor in volts (V), w is the angular speed
of the motor in radians per second (1/s), and L is the inductance due to the
windings in henries (H). The terms on the right-hand side are the mechanical
power produced by the motor, the power lost to heating the windings due to
the resistance of the wires, the power consumed or produced by energizing or
de-energizing the inductance of the windings (since the energy stored in an
inductor is 21 LI 2 , and power is the time derivative of energy), and the power


Chapter 8. Dynamics of Open Chains                                                                  307

                            Encoder
                                        Commutator
                                              Magnet
                                                         Windings
                                                                         Motor shaft

                                                                           Motor output pinion
                                                                           (input to gearhead)
   Torsional
   spring
                                                                                       Planetary
                                                                                       gearhead
     Brush        Housing
                  (magnetic return)

                                Ball bearings

                                                                                      Gearhead
      Commutator Windings                                                             output shaft

                                                   Magnet
                                                      Shaft

           Brush          Housing

\begin{figure}[h]
\centering
\includegraphics{figure_8_9.pdf}
\caption{(Top) A cutaway view of a Maxon brushed DC motor with an encoder}
\label{fig:8_9}
\end{figure}
and gearhead. (Cutaway image courtesy of Maxon Precision Motors, Inc., maxonmo-
torusa.com.) The motor’s rotor consists of the windings, commutator ring, and shaft.
Each of the several windings connects different segments of the commutator, and as
the motor rotates, the two brushes slide over the commutator ring and make contact
with different segments, sending current through one or more windings. One end of the
motor shaft turns the encoder, and the other end is input to the gearhead. (Bottom)
A simplified cross-section of the motor only, showing the stator (brushes, housing, and
magnets) in dark gray and the rotor (windings, commutator, and shaft) in light gray.

lost to friction in bearings, etc. Dropping this last term, replacing τ w by kt Iw,
and dividing by I, we get the voltage equation
                                                              dI
\begin{equation}
V = kt w + IR + L           .
\label{eq:8.102}
\end{equation}
                                                              dt


308                                                      8.9. Actuation, Gearing, and Friction

Often Equation (8.102) is written with the electrical constant ke (with units
of V s) instead of the torque constant kt , but in SI units (V s or N m/A) the
numerical values of the two are identical; they represent the same constant
property of the motor. So we prefer to use kt .
    The voltage term kt w in Equation (8.102) is called the back electromotive
force or back-emf for short, and it is what differentiates a motor from being
simply a resistor and inductor in series. It also allows a motor, which we usually
think of as converting electrical power to mechanical, to be run as a generator,
converting mechanical power to electrical. If the motor’s electrical inputs are
disconnected (so no current can flow) and the shaft is forced to turn by some
external torque, you can measure the back-emf voltage kt w across the motor’s
inputs.
    For simplicity, in the rest of this section we ignore the L dI/dt term. This as-
sumption is exactly satisfied when the motor is operating at a constant current.
With this assumption, Equation (8.102) can be rearranged to
                                       1             V    R
                                w=        (V − IR) =    − 2 τ,
                                       kt            kt  kt

expressing the speed w as a linear function of τ (with a slope of −R/kt2 ) for a con-
stant V . Now assume that the voltage across the motor is limited to the range
[−Vmax , +Vmax ] and the current through the motor is limited to [−Imax , +Imax ],
perhaps by the amplifier or power supply. Then the operating region of the mo-
tor in the torque–speed plane is as shown in Figure 8.10. Note that the signs of τ
and w are opposite in the second and fourth quadrants of this plane, and there-
fore the product τ w is negative. When the motor operates in these quadrants,
it is actually consuming mechanical power, not producing mechanical power.
The motor is acting like a damper.
    Focusing on the first quadrant (τ ≥ 0, w ≥ 0, τ w ≥ 0), the boundary of
the operating region is called the speed–torque curve. The no-load speed
w0 = Vmax /kt at one end of the speed–torque curve is the speed at which
the motor spins when it is powered by Vmax but is providing no torque. In
this operating condition, the back-emf kt w is equal to the applied voltage, so
there is no voltage remaining to create current (or torque). The stall torque
τstall = kt Vmax /R at the other end of the speed–torque curve is achieved when
the shaft is blocked from spinning, so there is no back-emf.
    Figure 8.10 also indicates the continuous operating region where |τ | ≤ τcont .
The motor may be operated intermittently outside the continuous operating
region, but extended operation outside the continuous operating region raises
the possibility that the motor will overheat.
    The motor’s rated mechanical power is Prated = τcont wcont , where wcont is


Chapter 8. Dynamics of Open Chains                                                                       309

                                             w
                       +V
                        ma
         −Imax limit

                             x l
                                   im
                                        it

                                                               rated mechanical power =
                                                  w0            τcont wcont

                                                       τcont    τstall                               τ

                                                                                       +Imax limit
                             continuous
                              operating
                               region                          −V
                                                                  ma
                                                                       x l
                                                                             im
                                                                                  it

\begin{figure}[h]
\centering
\includegraphics{figure_8_10.pdf}
\caption{The operating region (dark gray) of a current- and voltage-limited DC}
\label{fig:8_10}
\end{figure}
electric motor, and its continuous operating region (light gray).

the speed on the speed–torque curve corresponding to τcont . Even if the motor’s
rated power is sufficient for a particular application, the torque generated by
a DC motor is typically too low to be useful. As mentioned earlier, gearing is
therefore used to increase the torque while also decreasing the speed. For a gear
ratio G, the output speed of the gearhead is
                                                       wmotor
                                             wgear =          .
                                                        G
For an ideal gearhead, no power is lost in the torque conversion, so τmotor wmotor =
τgear wgear , which implies that
                                             τgear = Gτmotor .
In practice, some mechanical power is lost due to friction and impacts between
gear teeth, bearings, etc., so
                                             τgear = ηGτmotor ,
where η ≤ 1 is the efficiency of the gearhead.


310                                                      8.9. Actuation, Gearing, and Friction

                                           w

                                                                                         τ

                                                                           G=2

\begin{figure}[h]
\centering
\includegraphics{figure_8_11.pdf}
\caption{The original motor operating region, and the operating region with a}
\label{fig:8_11}
\end{figure}
gear ratio G = 2 showing the increased torque and decreased speed.

   Figure 8.11 shows the operating region of the motor from Figure 8.10 when
the motor is geared by G = 2 (with η = 1). The maximum torque doubles,
while the maximum speed shrinks by a factor of two. Since many DC motors
are capable of no-load speeds of 10 000 rpm or more, robot joints often have
gear ratios of 100 or more to achieve an appropriate compromise between speed
and torque.

\subsection{Apparent Inertia}
The motor’s stator is attached to one link and the rotor is attached to another
link, possibly through a gearhead. Therefore, when calculating the contribution
of a motor to the masses and inertias of the links, the mass and inertia of the
stator must be assigned to one link and the mass and inertia of the rotor must
be assigned to the other link.
    Consider a stationary link 0 with the stator of the joint-1 gearmotor attached
to it. The rotational speed of joint 1, the output of the gearhead, is θ̇. Therefore
the motor’s rotor rotates at Gθ̇. The kinetic energy of the rotor is therefore
                                 1                 1
                          K=       Irotor (Gθ̇)2 =          G2 Irotor      θ̇2 ,
                                 2                 2        | {z }
                                                        apparent inertia

where Irotor is the rotor’s scalar inertia about the rotation axis and G2 Irotor is
the apparent inertia (often called the reflected inertia) of the rotor about
the axis. In other words, if you were to grab link 1 and rotate it manually, the


Chapter 8. Dynamics of Open Chains                                                                  311

inertia contributed by the rotor would feel as if it were a factor G2 larger than
its actual inertia, owing to the gearhead.
    While the inertia Irotor is typically much less than the inertia Ilink of the
rest of the link about the rotation axis, the apparent inertia G2 Irotor may be
on the order of, or even larger than, Ilink .
    One consequence as the gear ratio becomes large is that the inertia seen by
joint i becomes increasingly dominated by the apparent inertia of the rotor. In
other words, the torque required of joint i becomes relatively more dependent
on θ̈i than on other joint accelerations, i.e., the robot’s mass matrix becomes
more diagonal. In the limit when the mass matrix has negligible off-diagonal
components (and in the absence of gravity), the dynamics of the robot are
decoupled – the dynamics at one joint has no dependence on the configuration
or motion of the other joints.
    As an example, consider the 2R arm of Figure 8.1 with L1 = L2 = m1 =
m2 = 1. Now assume that each of joint 1 and joint 2 has a motor of mass 1,
with a stator of inertia 0.005 and a rotor of inertia 0.00125, and a gear ratio G
(with η = 1). With a gear ratio G = 10, the mass matrix is
                                                            
                               4.13 + 2 cos θ2 1.01 + cos θ2
                    M (θ) =                                    .
                                1.01 + cos θ2       1.13

    With a gear ratio G = 100, the mass matrix is
                                                         
                            16.5 + 2 cos θ2 1.13 + cos θ2
                  M (θ) =                                   .
                             1.13 + cos θ2      13.5

The off-diagonal components are relatively less important for this second robot.
The available joint torques of the second robot are ten times that of the first
robot so, despite the increases in the mass matrix elements, the second robot is
capable of significantly higher accelerations and end-effector payloads. The top
speed of each joint of the second robot is ten times less than that of the first
robot, however.
    If the apparent inertia of the rotor is non-negligible relative to the inertia
of the rest of the link, the Newton–Euler inverse dynamics algorithm must be
modified to account for it. One approach is to treat the link as consisting of
two separate bodies, the geared rotor driving the link and the rest of the link,
each with its own center of mass and inertial properties (where the link inertial
properties include the inertial properties of the stator of any motor mounted on
the link). In the forward iteration, the twist and acceleration of each body is
determined while accounting for the gearhead in calculating the rotor’s motion.
In the backward iteration, the wrench on the link is calculated as the sum of two


312                                                      8.9. Actuation, Gearing, and Friction

                          joint i axis                                − AdTT(i+1) ,i (F(i+1)L )
                                                                                    L   L

                                                    link i
               rotor i axis
                                      gearhead
                                                                              {i + 1}R

                  rotor                                                                 {i}L
                 stator               motor i

                                       link i − 1                    FiL
                              (a)                                                (b)

\begin{figure}[h]
\centering
\includegraphics{figure_8_12.pdf}
\caption{(a) Schematic of a geared motor between links i − 1 and i. (b) The}
\label{fig:8_12}
\end{figure}
free-body diagram for link i, which is analogous to Figure 8.6.

wrenches: (i) the link wrench as given by Equation (8.53) and (ii) the reaction
wrench from the distal rotor. The resultant wrench projected onto the joint
axis is then the gear torque τgear ; dividing τgear by the gear ratio and adding
to this the torque resulting from the acceleration of the rotor results in the
required motor torque τmotor . The current command to the DC motor is then
Icom = τmotor /(ηkt ).

\subsection{Newton–Euler Inverse Dynamics Algorithm Accounting}
            for Motor Inertias and Gearing
We now reformulate the recursive Newton–Euler inverse dynamics algorithm
taking into account the apparent inertias as discussed above. Figure 8.12 il-
lustrates the setup. We assume massless gears and shafts and that the friction
between gears as well as the friction between shafts and links is negligible.

Initialization Attach a frame {0}L to the base, frames {1}L to {n}L to the
centers of mass of links 1 to n, and frames {1}R to {n}R to the centers of mass of
rotors 1 to n. Frame {n + 1}L is attached to the end-effector, which is assumed
fixed with respect to frame {n}L . Define MiR ,(i−1)L and MiL ,(i−1)L to be the
configuration of {i − 1}L in {i}R and in {i}L , respectively, when θi = 0. Let Ai
be the screw axis of joint i expressed in {i}L . Similarly, let Ri be the screw axis
of rotor i expressed in {i}R . Let GiL be the 6 × 6 spatial inertia matrix of link
i that includes the inertia of the attached stator and GiR be the 6 × 6 spatial


Chapter 8. Dynamics of Open Chains                                                                  313

inertia matrix of rotor i. The gear ratio of motor i is Gi . The twists V0L and
V̇0L and the wrench F(n+1)L are defined in the same way as V0 , V̇0 , and Fn+1
in Section 8.3.2.

Forward iterations Given θ, θ̇, θ̈, for i = 1 to n do
\begin{equation}
TiR ,(i−1)L     = e−[Ri ]Gi θi MiR ,(i−1)L ,
\label{eq:8.103}
\end{equation}
                                    −[Ai ]θi
\begin{equation}
TiL ,(i−1)L   = e                  MiL ,(i−1)L ,
\label{eq:8.104}
\end{equation}
\begin{equation}
ViR    =        AdTiR ,(i−1)L (V(i−1)L ) + Ri Gi θ̇i ,
\label{eq:8.105}
\end{equation}
\begin{equation}
ViL    =        AdTiL ,(i−1)L (V(i−1)L ) + Ai θ̇i ,
\label{eq:8.106}
\end{equation}
\begin{equation}
V̇iR   =        AdTiR ,(i−1)L (V̇(i−1)L ) + adViR (Ri )Gi θ̇i + Ri Gi θ̈i ,
\label{eq:8.107}
\end{equation}
\begin{equation}
V̇iL   =        AdTiL ,(i−1)L (V̇(i−1)L ) + adViL (Ai )θ̇i + Ai θ̈i .
\label{eq:8.108}
\end{equation}
Backward iterations For i = n to 1 do

    FiL       =    AdT                                   T
                     T(i+1) ,i (F(i+1)L ) + GiL V̇iL − adVi (GiL ViL )
                               L   L                                 L

                   + AdT                               T
                       T(i+1) ,i (G(i+1)R V̇(i+1)R − adV(i+1) (G(i+1)R V(i+1)R )),(8.109)
                                    R   L                                R
\begin{equation}
τi,gear     =    ATi FiL ,
\label{eq:8.110}
\end{equation}
                   τi,gear                    T
τi,motor      =            + RT
\begin{equation}
i (GiR V̇iR − adViR (GiR ViR )).
\label{eq:8.111}
\end{equation}
                     Gi
    In the backward iteration stage, the quantity F(n+1)L occurring in the first
step of the backward iteration is taken to be the external wrench applied to the
end-effector (expressed in the {n + 1}L frame), with G(n+1)R set to zero; FiL
denotes the wrench applied to link i via the motor i gearhead (expressed in the
{iL } frame); τi,gear is the torque generated at the motor i gearhead; and τi,motor
is the torque at rotor i.
    Note that if there is no gearing then no modification to the original Newton–
Euler inverse dynamics algorithm is necessary; the stator is attached to one link
and the rotor is attached to another link. Robots constructed with a motor at
each axis and no gearheads are sometimes called direct-drive robots. Direct-
drive robots have low friction, but they see limited use because typically the
motors must be large and heavy to generate appropriate torques.
    No modification is needed to the Lagrangian approach to the dynamics to
handle geared motors, provided that we can correctly represent the kinetic en-
ergy of the faster-spinning rotors.


314                                                      8.9. Actuation, Gearing, and Friction

      τfric          τfric            τfric           τfric            τfric           τfric

                θ̇               θ̇              θ̇              θ̇               θ̇              θ̇

         (a)             (b)              (c)             (d)             (e)              (f)

\begin{figure}[h]
\centering
\includegraphics{figure_8_13.pdf}
\caption{Examples of velocity-dependent friction models. (a) Viscous friction,}
\label{fig:8_13}
\end{figure}
τfric = bviscous θ̇. (b) Coulomb friction, τfric = bstatic sgn(θ̇); τfric can take any value in
[−bstatic , bstatic ] at zero velocity. (c) Static plus viscous friction, τfric = bstatic sgn(θ̇)+
bviscous θ̇. (d) Static and kinetic friction, requiring τfric ≥ |bstatic | to initiate motion and
then τfric = bkinetic sgn(θ̇) during motion, where bstatic > bkinetic . (e) Static, kinetic,
and viscous friction. (f) A friction law exhibiting the Stribeck effect – at low velocities,
the friction decreases as the velocity increases.

\subsection{Friction}
The Lagrangian and Newton–Euler dynamics do not account for friction at
the joints, but the friction forces and torques in gearheads and bearings may be
significant. Friction is a complex phenomenon that is the subject of considerable
current research; any friction model is a gross attempt to capture the average
behavior of the micromechanics of contact.
    Friction models often include a static friction term and a velocity-dependent
viscous friction term. The presence of a static friction term means that a
nonzero torque is required to cause the joint to begin to move. The viscous fric-
tion term indicates that the amount of friction torque increases with increasing
velocity of the joint. See Figure 8.13 for some examples of velocity-dependent
friction models.
    Other factors may contribute to the friction at a joint, including the loading
of the joint bearings, the time the joint has been at rest, the temperature, etc.
The friction in a gearhead often increases as the gear ratio G increases.

\subsection{Joint and Link Flexibility}
In practice, a robot’s joints and links are likely to exhibit some flexibility. For
example, the flexspline element of a harmonic drive gearhead achieves essentially
zero backlash by being somewhat flexible. A model of a joint with harmonic
drive gearing, then, could include a relatively stiff torsional spring between the
motor’s rotor and the link to which the gearhead is attached.
    Similarly, links themselves are not infinitely stiff. Their finite stiffness is
exhibited as vibrations along the link.


Chapter 8. Dynamics of Open Chains                                                                  315

    Flexible joints and links introduce extra states to the dynamics of the robot,
significantly complicating the dynamics and control. While many robots are
designed to be stiff in order to minimize these complexities, in some cases this
is impractical owing to the extra link mass required to create the stiffness.

\section{Summary}
    • Given a set of generalized coordinates θ and generalized forces τ , the
      Euler–Lagrange equations can be written
                                                  d ∂L ∂L
                                            τ=            −    ,
                                                  dt ∂ θ̇   ∂θ
       where L(θ, θ̇) = K(θ, θ̇) − P(θ), K is the kinetic energy of the robot, and
       P is the potential energy of the robot.
    • The equations of motion of a robot can be written in the following equiv-
      alent forms:
                                  τ    =    M (θ)θ̈ + h(θ, θ̇)
                                       =    M (θ)θ̈ + c(θ, θ̇) + g(θ)
                                       =    M (θ)θ̈ + θ̇T Γ(θ)θ̇ + g(θ)
                                       =    M (θ)θ̈ + C(θ, θ̇)θ̇ + g(θ),
       where M (θ) is the n×n symmetric positive-definite mass matrix, h(θ, θ̇) is
       the sum of the generalized forces due to the gravity and quadratic velocity
       terms, c(θ, θ̇) are quadratic velocity forces, g(θ) are gravitational forces,
       Γ(θ) is an n×n×n matrix of Christoffel symbols of the first kind obtained
       from partial derivatives of M (θ) with respect to θ, and C(θ, θ̇) is the n × n
       Coriolis matrix whose (i, j)th entry is given by
                                                        n
                                                        X
                                        cij (θ, θ̇) =         Γijk (θ)θ̇k .
                                                        k=1

       If the end-effector of the robot is applying a wrench Ftip to the environ-
       ment, the term J T (θ)Ftip should be added to the right-hand side of the
       robot’s dynamic equations.
    • The symmetric positive-definite rotational inertia matrix of a rigid body
      is                                            
                                     Ixx Ixy Ixz
                            Ib =  Ixy Iyy Iyz  ,
                                     Ixz Iyz Izz


316                                                                                  8.10. Summary

        where
                   R                                             R
             Ixx = R B (y 2 + z 2 )ρ(x, y, z)dV,           Iyy = BR(x2 + z 2 )ρ(x, y, z)dV,
             Izz = BR(x2 + y 2 )ρ(x, y, z)dV,              Ixy = − R B xyρ(x, y, z)dV,
             Ixz = − B xzρ(x, y, z)dV,                     Iyz = − B yzρ(x, y, z)dV,
        B is the volume of the body, dV is a differential volume element, and
        ρ(x, y, z) is the density function.
      • If Ib is defined in a frame {b} at the center of mass, with axes aligned
        with the principal axes of inertia, then Ib is diagonal.
      • If {b} is at the center of mass but its axes are not aligned with the principal
        axes of inertia, there always exists a rotated frame {c} defined by the
                                                 T
        rotation matrix Rbc such that Ic = Rbc     Ib Rbc is diagonal.
      • If Ib is defined in a frame {b} at the center of mass then Iq , the inertia
        in a frame {q} aligned with {b} but displaced from the origin of {b} by
        q ∈ R3 in {b} coordinates, is
                                       Iq = Ib + m(q T qI − qq T ).

      • The spatial inertia matrix Gb expressed in a frame {b} at the center of
        mass is defined as the 6 × 6 matrix
                                                  
                                           Ib 0
                                    Gb =             .
                                            0 mI

        In a frame {a} at a configuration Tba relative to {b}, the spatial inertia
        matrix is
                               Ga = [AdTba ]T Gb [AdTba ].

      • The Lie bracket of two twists V1 and V2 is
                                          adV1 (V2 ) = [adV1 ]V2 ,
        where                                                 
                                                   [ω]    0
                                    [adV ] =                       ∈ R6×6 .
                                                   [v]   [ω]
      • The twist–wrench formulation of the rigid-body dynamics of a single rigid
        body is
                              Fb = Gb V̇b − [adVb ]T Gb Vb .
        The equations have the same form if F, V, and G are all expressed in the
        same frame, regardless of the frame.


Chapter 8. Dynamics of Open Chains                                                                  317

    • The kinetic energy of a rigid body is 12 VbT Gb Vb , and the kinetic energy of
      an open-chain robot is 12 θ̇T M (θ)θ̇.
    • The forward–backward Newton–Euler inverse dynamics algorithm is the
      following:
      Initialization: Attach a frame {0} to the base, frames {1} to {n} to the
      centers of mass of links {1} to {n}, and a frame {n+1} at the end-effector,
      fixed in the frame {n}. Define Mi,i−1 to be the configuration of {i − 1}
      in {i} when θi = 0. Let Ai be the screw axis of joint i expressed in {i},
      and Gi be the 6 × 6 spatial inertia matrix of link i. Define V0 to be the
      twist of the base frame {0} expressed in base-frame coordinates. (This
      quantity is typically zero.) Let g ∈ R3 be the gravity vector expressed in
      base-frame-{0} coordinates, and define V̇0 = (0, −g). (Gravity is treated
      as an acceleration of the base in the opposite direction.) Define Fn+1 =
      Ftip = (mtip , ftip ) to be the wrench applied to the environment by the
      end-effector expressed in the end-effector frame {n + 1}.
       Forward iterations: Given θ, θ̇, θ̈, for i = 1 to n do

                        Ti,i−1      = e−[Ai ]θi Mi,i−1 ,
                             Vi     =     AdTi,i−1 (Vi−1 ) + Ai θ̇i ,
                             V̇i    =     AdTi,i−1 (V̇i−1 ) + adVi (Ai )θ̇i + Ai θ̈i .

       Backward iterations: For i = n to 1 do

                          Fi       =    AdT                           T
                                          Ti+1,i (Fi+1 ) + Gi V̇i − adVi (Gi Vi ),

                           τi      = FiT Ai .

    • Let Jib (θ) be the Jacobian relating θ̇ to the body twist Vi in link i’s center-
      of-mass frame {i}. Then the mass matrix M (θ) of the manipulator can
      be expressed as
                                        Xn
                                               T
                               M (θ) =       Jib (θ)Gi Jib (θ).
                                                  i=1

    • The forward dynamics problem involves solving

                                   M (θ)θ̈ = τ (t) − h(θ, θ̇) − J T (θ)Ftip

       for θ̈, using any efficient solver of equations of the form Ax = b.


318                                                                                  8.10. Summary

      • The robot’s dynamics M (θ)θ̈ + h(θ, θ̇) can be expressed in the task space
        as
                                F = Λ(θ)V̇ + η(θ, V),
        where F is the wrench applied to the end-effector, V is the twist of the end-
        effector, and F, V, and the Jacobian J(θ) are all defined in the same frame.
        The task-space mass matrix Λ(θ) and gravity and quadratic velocity forces
        η(θ, V) are

                             Λ(θ)      = J −T M (θ)J −1 ,
                           η(θ, V)                                ˙ −1 V.
                                       = J −T h(θ, J −1 V) − Λ(θ)JJ

      • Define two n × n projection matrices of rank n − k

                     P (θ) = I − AT (AM −1 AT )−1 AM −1 ,
                     Pθ̈ (θ) = M −1 P M = I − M −1 AT (AM −1 AT )−1 A,

        corresponding to the k Pfaffian constraints, A(θ)θ̇ = 0, A ∈ Rk×n , acting
        on the robot. Then the n + k constrained equations of motion

                                  τ      = M (θ)θ̈ + h(θ, θ̇) + AT (θ)λ,
                              A(θ)θ̇     = 0

        can be reduced to the following equivalent forms by eliminating the La-
        grange multipliers λ:

                                       Pτ       = P (M θ̈ + h),
                                       Pθ̈ θ̈   = Pθ̈ M −1 (τ − h).

        The matrix P projects away the joint force–torque components that act
        on the constraints without doing work on the robot, and the matrix Pθ̈
        projects away acceleration components that do not satisfy the constraints.
      • An ideal gearhead (one that is 100% efficient) with a gear ratio G multi-
        plies the torque at the output of a motor by a factor G and divides the
        speed by the factor G, leaving the mechanical power unchanged. The in-
        ertia of the motor’s rotor about its axis of rotation, as it appears at the
        output of the gearhead, is G2 Irotor .


Chapter 8. Dynamics of Open Chains                                                                  319

\section{Software}
Software functions associated with this chapter are listed below.
adV = ad(V)
Computes [adV ].
taulist = InverseDynamics(thetalist,dthetalist,ddthetalist,g,Ftip,
Mlist,Glist,Slist)
Uses Newton–Euler inverse dynamics to compute the n-vector τ of the required
joint forces–torques given θ, θ̇, θ̈, g, Ftip , a list of transforms Mi−1,i specifying
the configuration of the center-of-mass frame of link {i} relative to {i − 1} when
the robot is at its home position, a list of link spatial inertia matrices Gi , and a
list of joint screw axes Si expressed in the base frame.
M = MassMatrix(thetalist,Mlist,Glist,Slist)
Computes the mass matrix M (θ) given the joint configuration θ, a list of trans-
forms Mi−1,i , a list of link spatial inertia matrices Gi , and a list of joint screw
axes Si expressed in the base frame.
c = VelQuadraticForces(thetalist,dthetalist,Mlist,Glist,Slist)
Computes c(θ, θ̇) given the joint configuration θ, the joint velocities θ̇, a list of
transforms Mi−1,i , a list of link spatial inertia matrices Gi , and a list of joint
screw axes Si expressed in the base frame.
grav = GravityForces(thetalist,g,Mlist,Glist,Slist)
Computes g(θ) given the joint configuration θ, the gravity vector g, a list of
transforms Mi−1,i , a list of link spatial inertia matrices Gi , and a list of joint
screw axes Si expressed in the base frame.
JTFtip = EndEffectorForces(thetalist,Ftip,Mlist,Glist,Slist)
Computes J T (θ)Ftip given the joint configuration θ, the wrench Ftip applied by
the end-effector, a list of transforms Mi−1,i , a list of link spatial inertia matrices
Gi , and a list of joint screw axes Si expressed in the base frame.
ddthetalist = ForwardDynamics(thetalist,dthetalist,taulist,g,Ftip,
Mlist,Glist,Slist)
Computes θ̈ given the joint configuration θ, the joint velocities θ̇, the joint forces-
torques τ , the gravity vector g, the wrench Ftip applied by the end-effector, a
list of transforms Mi−1,i , a list of link spatial inertia matrices Gi , and a list of
joint screw axes Si expressed in the base frame.
[thetalistNext,dthetalistNext] = EulerStep(thetalist,dthetalist,


320                                                                   8.12. Notes and References

ddthetalist,dt)
Computes a first-order Euler approximation to {θ(t + δt), θ̇(t + δt)} given the
joint configuration θ(t), the joint velocities θ̇(t), the joint accelerations θ̈(t), and
a timestep δt.
taumat = InverseDynamicsTrajectory(thetamat,dthetamat,ddthetamat,
g,Ftipmat,Mlist,Glist,Slist)
The variable thetamat is an N × n matrix of robot joint variables θ, where the
ith row corresponds to the n-vector of joint variables θ(t) at time t = (i − 1)δt,
where δt is the timestep. The variables dthetamat, ddthetamat, and Ftipmat
similarly represent θ̇, θ̈, and Ftip as a function of time. Other inputs include the
gravity vector g, a list of transforms Mi−1,i , a list of link spatial inertia matrices
Gi , and a list of joint screw axes Si expressed in the base frame. This function
computes an N × n matrix taumat representing the joint forces-torques τ (t)
required to generate the trajectory specified by θ(t) and Ftip (t). Note that it is
not necessary to specify δt. The velocities θ̇(t) and accelerations θ̈(t) should be
consistent with θ(t).
[thetamat,dthetamat] = ForwardDynamicsTrajectory(thetalist,
dthetalist,taumat,g,Ftipmat,Mlist,Glist,Slist,dt,intRes)
This function numerically integrates the robot’s equations of motion using Euler
integration. The outputs are N × n matrices thetamat and dthetamat, where
the ith rows correspond respectively to the n-vectors θ((i−1)δt) and θ̇((i−1)δt).
The inputs are the initial state θ(0), θ̇(0), an N × n matrix of joint forces or
torques τ (t), the gravity vector g, an N × n matrix of end-effector wrenches
Ftip (t), a list of transforms Mi−1,i , a list of link spatial inertia matrices Gi , a
list of joint screw axes Si expressed in the base frame, the timestep δt, and the
number of integration steps to take during each timestep (a positive integer).

\section{Notes and References}
An accessible general reference on rigid-body dynamics that covers both the
Newton–Euler and Lagrangian formulations is [51]. A more classical reference
that covers a wide range of topics in dynamics is [193].
    A recursive inverse dynamics algorithm for open chains using the classical
screw-theoretic machinery of twists and wrenches was first formulated by Feath-
erstone (the collection of twists, wrenches, and the corresponding analogues of
accelerations, momentum, and inertias, are collectively referred to as spatial
vector notation); this formulation, as well as more efficient extensions based on
articulated body inertias, are described in [46, 47].


Chapter 8. Dynamics of Open Chains                                                                  321

    The recursive inverse dynamics algorithm presented in this chapter was first
described in [133] and makes use of standard operators from the theory of Lie
groups and Lie algebras. One of the important practical advantages of this ap-
proach is that analytic formulas can be derived for taking first- and higher-order
derivatives of the dynamics. This has important consequences for dynamics-
based motion optimization: the availability of analytic gradients can greatly
improve the convergence and robustness for motion optimization algorithms.
These and other related issues are explored in, e.g., [88].
    The task-space formulation was first initiated by Khatib [74], who referred
to it as the operational space formulation. Note that the task-space formula-
tion involves taking time derivatives of the forward kinematics Jacobian, i.e.,
 ˙
J(θ).  Using either the body or space Jacobian, J(θ) ˙    can in fact be evaluated
analytically; this is explored in one of the exercises at the end this chapter.
    A brief history of the evolution of robot dynamics algorithms, as well as
pointers to references on the more general subject of multibody system dynamics
(of which robot dynamics can be considered a subfield) can be found in [48].

\section{Exercises}

Exercise 8.1 Derive the formulas given in Figure 8.5 for:
 (a) a rectangular parallelpiped;
 (b) a circular cylinder;
 (c) an ellipsoid.

Exercise 8.2 Consider a cast iron dumbbell consisting of a cylinder connecting
two solid spheres at either end of the cylinder. The density of the dumbbell is
7500 kg/m3 . The cylinder has a diameter of 4 cm and a length of 20 cm. Each
sphere has a diameter of 20 cm.
  (a) Find the approximate rotational inertia matrix Ib in a frame {b} at the
      center of mass with axes aligned with the principal axes of inertia of the
      dumbbell.
 (b) Write down the spatial inertia matrix Gb .

Exercise 8.3 Rigid-body dynamics in an arbitrary frame.
 (a) Show that Equation (8.42) is a generalization of Steiner’s theorem.
 (b) Derive Equation (8.43).

Exercise 8.4 The 2R open-chain robot of Figure 8.14, referred to as a ro-


322                                                                                        8.13. Exercises

                            θ1
                            τ1       ẑs                                              θ1
                    L1                                                                τ1
                                               ŷs
                               x̂s                                          L1
             ẑ1                     {s}                                                   ẑs
        {b1 }
                                                                 {b } ẑ1
                                                                    1            m1
                                                                                                 ŷs
          x̂     ŷ1                                                                   x̂s {s}
       τ2 1                          rod 1
      θ2                  L2                                        τ x̂1    ŷ1
                                                                  θ2 2
              ẑ2   ŷ2                    g                                     L2 g
                                                                         ẑ2
                                                                             ŷ2
            x̂2 {b2 }                                                  x̂2
                                     Zero Position                                 Zero Position
                 rod 2                                             {b2 } m2

\begin{figure}[h]
\centering
\includegraphics{figure_8_14.pdf}
\caption{2R rotational inverted pendulum. (Left) Its construction; (right) the}
\label{fig:8_14}
\end{figure}
model.

tational inverted pendulum or Furuta pendulum, is shown in its zero position.
Assuming that the mass of each link is concentrated at the tip and neglecting
its thickness, the robot can be modeled as shown in the right-hand figure. As-
sume that m1 = m2 = 2, L1 = L2 = 1, g = 10, and the link inertias I1 and I2
(expressed in their respective link frames {b1 } and {b2 }) are
                                                          
                          0 0 0                    4 0 0
                   I1 =  0 4 0  ,        I2 =  0 4 0  .
                          0 0 4                    0 0 0
  (a) Derive the dynamic equations and determine the input torques τ1 and τ2
      when θ1 = θ2 = π/4 and the joint velocities and accelerations are all zero.
  (b) Draw the torque ellipsoid for the mass matrix M (θ) when θ1 = θ2 = π/4.

Exercise 8.5 Prove the following Lie bracket identity (called the Jacobi iden-
tity) for arbitrary twists V1 , V2 , V3 :
               adV1 (adV2 (V3 )) + adV3 (adV1 (V2 )) + adV2 (adV3 (V1 )) = 0.

                                  ˙
Exercise 8.6 The evaluation of J(θ),    the time derivative of the forward kine-
matics Jacobian, is needed in the calculation of the frame accelerations V̇i in
the Newton–Euler inverse dynamics algorithm and also in the formulation of
the task-space dynamics. Letting Ji (θ) denote the ith column of J(θ), we have
                                                           Xn
                                               d               ∂Ji
                                                  Ji (θ) =         θ̇j .
                                               dt          j=1
                                                               ∂θj


Chapter 8. Dynamics of Open Chains                                                                  323

  (a) Suppose that J(θ) is the space Jacobian. Show that
                                  
                           ∂Ji      adJi (Jj ) for i > j
                               =
                           ∂θj              0 for i ≤ j.
  (b) Now suppose that J(θ) is the body Jacobian. Show that
                                 
                         ∂Ji        adJi (Jj ) for i < j
                              =
                         ∂θj                0 for i ≥ j.

Exercise 8.7 Show that the time derivative of the mass matrix Ṁ (θ) can be
written explicitly as
                 Ṁ = AT LT ΓT [adAθ̇ ]T LT GLA + AT LT GL[adAθ̇ ]ΓA,
with the matrices as defined in the closed-form dynamics formulation.

Exercise 8.8 Explain intuitively the shapes of the end-effector force ellipsoids
in Figure 8.4 on the basis of the point masses and the Jacobians.

Exercise 8.9 Consider a motor with rotor inertia Irotor connected through
a gearhead of gear ratio G to a load with scalar inertia Ilink about the rota-
tion axis. The load and motor are said to be inertia matched if, for any
given torque τm at the motor, the acceleration of the load is maximized. The
acceleration of the load can be written
                                         Gτm
                              θ̈ =                   .
                                   Ilink + G2 Irotor
                                           p
Solve for the inertia-matching gear ratio Ilink /Irotor by solving dθ̈/dG = 0.

Exercise 8.10 Give the steps that rearrange Equation (8.98) to get Equa-
tion (8.100). Remember that P (θ) is not full rank and cannot be inverted.

Exercise 8.11 Program a function to calculate h(θ, θ̇) = c(θ, θ̇) + g(θ) effi-
ciently using Newton–Euler inverse dynamics.

Exercise 8.12 Give the equations that would convert the joint and link de-
scriptions in a robot’s URDF file to the data Mlist, Glist, and Slist, suitable
for using with the Newton–Euler algorithm InverseDynamicsTrajectory.

Exercise 8.13 The efficient evaluation of M (θ).


324                                                                                  8.13. Exercises

  (a) Develop conceptually a computationally efficient algorithm for determin-
      ing the mass matrix M (θ) using Equation (8.57).
  (b) Implement this algorithm.

Exercise 8.14 The function InverseDynamicsTrajectory requires the user
to enter not only a time sequence of joint variables thetamat but also a time
sequence of joint velocities dthetamat and accelerations ddthetamat. Instead,
the function could use numerical differencing to find approximately the joint
velocities and accelerations at each timestep, using only thetamat. Write an
alternative InverseDynamicsTrajectory function that does not require the
user to enter dthetamat and ddthetamat. Verify that it yields similar results.

Exercise 8.15 Dynamics of the UR5 robot.
 (a) Write the spatial inertia matrices Gi of the six links of the UR5, given
     the center-of-mass frames and mass and inertial properties defined in the
     URDF in Section 4.2.
 (b) Simulate the UR5 falling under gravity with acceleration g = 9.81 m/s2
     in the −ẑs -direction. The robot starts at its zero configuration and zero
     joint torques are applied. Simulate the motion for three seconds, with at
     least 100 integration steps per second. (Ignore the effects of friction and
     the geared rotors.)



\end{document}
